%%%%%%%%%%%%%%%%%%%%%%%%%%%%%%%%%%%%%%%%%%%%%%%%%%%%%%%%%%%%%%%%%%%%%%%%%%%%%%%
%% Trabalho de Conclusão de Curso --- Engenharia de Software, iCEV
%%
%% Classe: abnTeX2 (ABNT NBR 14724, 6023, 6024, 6027, 6028, 6034, 10520).
%% Substitui o modelo POLI/UFRJ usado no rascunho inicial, cuja classe fixava
%% nome, logotipo, grau e cidade da UFRJ em oito pontos da capa, folha de rosto
%% e resumo, e não admitia o curso de Engenharia de Software.
%%
%% COMPILAR:   make          (ou: make limpo && make)
%% O Makefile executa pdflatex -> bibtex -> makeglossaries -> makeindex ->
%% pdflatex x2, na ordem verificada.
%%
%% Este diretório vive dentro do repositório, em docs/assembly_tcc_mount/, e
%% consome do próprio repositório, sem copiar nada:
%%   figuras       ../../experiments/figures/output/   (regeneráveis)
%%   hardware      ../../hardware/
%%   referências   ../relatorio/referencias.bib          (fonte canônica)
%%%%%%%%%%%%%%%%%%%%%%%%%%%%%%%%%%%%%%%%%%%%%%%%%%%%%%%%%%%%%%%%%%%%%%%%%%%%%%%

\documentclass[
  12pt,
  openright,
  oneside,
  a4paper,
  english,
  brazil
]{abntex2}

%% --- Fontes e codificação --------------------------------------------------
\usepackage[T1]{fontenc}
\usepackage[utf8]{inputenc}
%% Times no corpo. As figuras do pipeline (experiments/figures/scripts/
%% theme_kaelix.R) compõem em Times New Roman para casar com este corpo;
%% trocar a fonte aqui sem trocar lá reabre a incoerência visual.
\usepackage{newtxtext,newtxmath}
\usepackage{microtype}

%% --- Pacotes de conteúdo ---------------------------------------------------
\usepackage{graphicx}
\usepackage{xcolor}      % marcas vermelhas \emandamento
\usepackage{booktabs}    % réguas de tabela
\usepackage{amsmath}
\usepackage{float}
\usepackage{caption}
\usepackage{etoolbox}    % \ifstrequal em simbolos.tex
\usepackage{siunitx}
%% Separador decimal do português: vírgula. Sem isto, \num{0,5} sai "0.5".
\sisetup{
  output-decimal-marker = {,},
  group-separator       = {\,},
  range-phrase          = {~a~},
  list-final-separator  = {~e~},
  per-mode              = symbol
}

%% --- Citações (ABNT, sistema autor-data) -----------------------------------
\usepackage[alf]{abntex2cite}

%% --- Listas de siglas e de símbolos ----------------------------------------
%% glossaries deve vir DEPOIS do hyperref, que a abntex2 carrega.
\usepackage[acronym,nomain,toc,nonumberlist]{glossaries}
\makeglossaries
\usepackage[intoc]{nomencl}
\makenomenclature

%% --- Caminhos para o repositório -------------------------------------------
%% Definidos uma única vez aqui; os capítulos têm fallback idêntico.
\newcommand{\kxfig}[1]{../../experiments/figures/output/#1}
\newcommand{\kxhw}[1]{../../#1}

%%%%%%%%%%%%%%%%%%%%%%%%%%%%%%%%%%%%%%%%%%%%%%%%%%%%%%%%%%%%%%%%%%%%%%%%%%%%%%%
%% DADOS DO TRABALHO
%% Campos entre colchetes ainda precisam ser preenchidos.
%%%%%%%%%%%%%%%%%%%%%%%%%%%%%%%%%%%%%%%%%%%%%%%%%%%%%%%%%%%%%%%%%%%%%%%%%%%%%%%

%% TÍTULO PROVISÓRIO (working title). Finalizar só após o ensaio do protótipo:
%% o título deve refletir o que foi realizado, e os resultados ainda podem mudar.
%%   - Descritivo, sem afirmar resultado ("limites" entra só depois do ensaio).
%%   - Sem nome de projeto (Kaelix) nem modelo comercial (MPU6050, ESP32-S3):
%%     a NBR 14724 exige que o título permita indexação e recuperação.
%%   - Sem sigla (MEMS): desaconselhada em título por IEEE e Nature.
%%   - O subtítulo descreve o conteúdo; na versão final, é onde entra a
%%     contribuição apurada.
\titulo{Monitoramento de vibração de baixo custo para manutenção preditiva: aquisição e detecção embarcada de anomalias em motores}
\autor{Lauan Matheus}
\local{Teresina -- PI}
\data{2026}
\orientador{[Nome do orientador]}
%\coorientador{[Nome do coorientador]}
\instituicao{%
  Instituto de Ensino Superior iCEV
  \par
  Bacharelado em Engenharia de Software}
\tipotrabalho{Trabalho de Conclusão de Curso}
\preambulo{Trabalho de Conclusão de Curso apresentado ao curso de Bacharelado em
Engenharia de Software do Instituto de Ensino Superior iCEV, como requisito parcial
para a obtenção do título de Bacharel em Engenharia de Software.}

%% Metadados do PDF
\makeatletter
\hypersetup{
  pdftitle    = {\@title},
  pdfauthor   = {\@author},
  pdfsubject  = {\imprimirpreambulo},
  pdfcreator  = {LaTeX com abnTeX2},
  pdfkeywords = {manutenção preditiva}{análise de vibração}{detecção de anomalia}
                {sistemas embarcados}{instrumentação de baixo custo},
  colorlinks  = true,
  linkcolor   = black,
  citecolor   = black,
  filecolor   = black,
  urlcolor    = black,
  bookmarksdepth = 4
}
\makeatother

%% Capa. A capa padrão da abnTeX2 omite a instituição; a NBR 14724 a prevê como
%% primeiro elemento da capa "quando solicitado", e a maioria das instituições
%% solicita. Remover o bloco da instituição se o iCEV dispensar.
\renewcommand{\imprimircapa}{%
  \begin{capa}%
    \center
    %% Caixa-alta escrita à mão: \MakeUppercase transformaria "iCEV" em "ICEV".
    {\ABNTEXchapterfont\large INSTITUTO DE ENSINO SUPERIOR iCEV}\par
    {\ABNTEXchapterfont\large Bacharelado em Engenharia de Software}
    \vspace*{2.5cm}

    {\ABNTEXchapterfont\large\imprimirautor}

    \vfill
    \begin{center}
      \ABNTEXchapterfont\bfseries\LARGE\imprimirtitulo
    \end{center}
    \vfill

    \large\imprimirlocal
    \par
    \large\imprimirdata
    \vspace*{1cm}
  \end{capa}%
}

%% Títulos em serifa. A abnTeX2 usa sans nos títulos por padrão, o que reabriria
%% a incoerência corrigida nas figuras: corpo, títulos e figuras em Times.
\renewcommand{\ABNTEXchapterfont}{\rmfamily}

%% Espaçamento de parágrafo (ABNT: recuo, sem espaço entre parágrafos)
\setlength{\parindent}{1.3cm}
\setlength{\parskip}{0.2cm}

\makeindex

%% Listas de siglas e símbolos (definidas antes do documento)
%%%%%%%%%%%%%%%%%%%%%%%%%%%%%%%%%%%%%%%%%%%%%%%%%%%%%%%%%%%%%%%%%%%%%%%%%%%%%%%
%% Lista de Siglas e Abreviaturas (pacote glossaries)
%%
%% Contém apenas as siglas efetivamente usadas nos Capítulos 1 a 5 e nos
%% apêndices. Acrescentar entrada aqui ao introduzir sigla nova no texto.
%%%%%%%%%%%%%%%%%%%%%%%%%%%%%%%%%%%%%%%%%%%%%%%%%%%%%%%%%%%%%%%%%%%%%%%%%%%%%%%

%% Instrumentação e mecânica
\newacronym{MEMS}{MEMS}{sistemas microeletromecânicos (\emph{Micro-Electro-Mechanical Systems})}
\newacronym{IMU}{IMU}{unidade de medição inercial (\emph{Inertial Measurement Unit})}
\newacronym{BPFO}{BPFO}{frequência de passagem de esferas na pista externa (\emph{Ball Pass Frequency, Outer race})}
\newacronym{BPFI}{BPFI}{frequência de passagem de esferas na pista interna (\emph{Ball Pass Frequency, Inner race})}
\newacronym{DLPF}{DLPF}{filtro passa-baixas digital (\emph{Digital Low-Pass Filter})}
\newacronym{RMS}{RMS}{valor eficaz (\emph{Root Mean Square})}
\newacronym{ASA}{ASA}{acrilonitrila estireno acrilato}
\newacronym{NTC}{NTC}{termistor de coeficiente de temperatura negativo}

%% Normalização
\newacronym{ISO}{ISO}{Organização Internacional de Normalização (\emph{International Organization for Standardization})}
\newacronym{ABNT}{ABNT}{Associação Brasileira de Normas Técnicas}

%% Detecção e avaliação
\newacronym{ROC}{ROC}{característica de operação do receptor (\emph{Receiver Operating Characteristic})}
\newacronym{AUC}{AUC}{área sob a curva (\emph{Area Under the Curve})}
\newacronym{pAUC}{pAUC}{área parcial sob a curva, restrita a uma faixa de falso alarme}
\newacronym{PRAUC}{PR-AUC}{área sob a curva de precisão e revocação}
\newacronym{TPR}{TPR}{taxa de detecção (\emph{True Positive Rate})}
\newacronym{DFT}{DFT}{transformada discreta de Fourier (\emph{Discrete Fourier Transform})}
\newacronym{FFT}{FFT}{transformada rápida de Fourier (\emph{Fast Fourier Transform})}
\newacronym{FPR}{FPR}{taxa de falso alarme (\emph{False Positive Rate})}

%% Computação embarcada e comunicação
\newacronym{TinyML}{TinyML}{aprendizado de máquina em dispositivos de recursos limitados}
\newacronym{LoRa}{LoRa}{modulação de longo alcance e baixa potência (\emph{Long Range})}
\newacronym{SF}{SF}{fator de espalhamento (\emph{Spreading Factor})}
\newacronym{CRC}{CRC}{verificação cíclica de redundância (\emph{Cyclic Redundancy Check})}
\newacronym{ADC}{ADC}{conversor analógico-digital}
\newacronym{SRAM}{SRAM}{memória estática de acesso aleatório (\emph{Static Random-Access Memory})}

%% Projeto de hardware
\newacronym{PCB}{PCB}{placa de circuito impresso (\emph{Printed Circuit Board})}
\newacronym{ERC}{ERC}{verificação de regras elétricas (\emph{Electrical Rules Check})}
\newacronym{DRC}{DRC}{verificação de regras de projeto (\emph{Design Rules Check})}
\newacronym{PDN}{PDN}{rede de distribuição de alimentação (\emph{Power Delivery Network})}

%% Instituição
\newacronym{UFRJ}{UFRJ}{Universidade Federal do Rio de Janeiro}

%%%%%%%%%%%%%%%%%%%%%%%%%%%%%%%%%%%%%%%%%%%%%%%%%%%%%%%%%%%%%%%%%%%%%%%%%%%%%%%
%% O pacote glossaries só imprime as siglas EFETIVAMENTE USADAS via \gls{...}.
%% Como o texto dos capítulos usa as siglas em forma literal (e não \gls{}),
%% \glsaddall força todas as entradas acima para a lista impressa.
%%
%% Se no futuro o texto passar a usar \gls{}, remova a linha abaixo para que a
%% lista volte a conter apenas o que é citado.
%%%%%%%%%%%%%%%%%%%%%%%%%%%%%%%%%%%%%%%%%%%%%%%%%%%%%%%%%%%%%%%%%%%%%%%%%%%%%%%
\glsaddall

%%%%%%%%%%%%%%%%%%%%%%%%%%%%%%%%%%%%%%%%%%%%%%%%%%%%%%%%%%%%%%%%%%%%%%%%%%%%%%%
%% Lista de Símbolos (pacote nomencl)
%%
%% Contém apenas os símbolos que aparecem nas equações dos Capítulos 3 e 4.
%% Grupos: [S] modelo de sinal, [M] modelo mecânico, [D] descritores e
%% avaliação. A macro \nomunit acrescenta a unidade alinhada à direita.
%%%%%%%%%%%%%%%%%%%%%%%%%%%%%%%%%%%%%%%%%%%%%%%%%%%%%%%%%%%%%%%%%%%%%%%%%%%%%%%

\renewcommand{\nomname}{Lista de Símbolos}

%----------------------------------------------
\newcommand{\nomunit}[1]{%
\renewcommand{\nomentryend}{\hspace*{\fill}#1}}
%----------------------------------------------
\renewcommand\nomgroup[1]{%
   \item[\bfseries
   \ifstrequal{#1}{S}{Modelo de sinal e aquisição}{%
   \ifstrequal{#1}{M}{Modelo mecânico e térmico}{%
   \ifstrequal{#1}{D}{Descritores e avaliação}{}}}%
]}

%% --- Modelo de sinal e aquisição -------------------------------------------
\nomenclature[S]{\(a(t)\)}{Aceleração simulada\nomunit{\si{\meter\per\second\squared}}}
\nomenclature[S]{\(f_r\)}{Frequência de rotação do eixo\nomunit{\si{\hertz}}}
\nomenclature[S]{\(A_k\)}{Amplitude da \(k\)-ésima harmônica da rotação\nomunit{\si{\meter\per\second\squared}}}
\nomenclature[S]{\(\varphi_k\)}{Fase da \(k\)-ésima harmônica\nomunit{\si{\radian}}}
\nomenclature[S]{\(D\)}{Amplitude do impacto do defeito\nomunit{\si{\meter\per\second\squared}}}
\nomenclature[S]{\(\zeta_d\)}{Taxa de amortecimento do impacto\nomunit{\si{\per\second}}}
\nomenclature[S]{\(f_{\mathrm{res}}\)}{Frequência da ressonância estrutural excitada\nomunit{\si{\hertz}}}
\nomenclature[S]{\(t_i\)}{Instante do \(i\)-ésimo impacto\nomunit{\si{\second}}}
\nomenclature[S]{\(\delta_i\)}{Desvio aleatório que modela o escorregamento das esferas\nomunit{\si{\second}}}
\nomenclature[S]{\(\varepsilon(t)\)}{Ruído gaussiano branco\nomunit{\si{\meter\per\second\squared}}}
\nomenclature[S]{\(u(\cdot)\)}{Degrau unitário}
\nomenclature[S]{\(f_s\)}{Taxa de amostragem\nomunit{\si{\hertz}}}
\nomenclature[S]{\(f_{\mathrm{ap}}\)}{Frequência aparente sob subamostragem sem filtro\nomunit{\si{\hertz}}}
\nomenclature[S]{\(v_{\mathrm{rms}}\)}{Velocidade eficaz de banda larga, indicador normativo\nomunit{\si{\milli\meter\per\second}}}

\nomenclature[S]{\(N_b\)}{Número de esferas do rolamento}
\nomenclature[S]{\(d\)}{Diâmetro da esfera\nomunit{\si{\milli\meter}}}
\nomenclature[S]{\(D_p\)}{Diâmetro primitivo do rolamento\nomunit{\si{\milli\meter}}}
\nomenclature[S]{\(\theta\)}{Ângulo de contato do rolamento\nomunit{\si{\radian}}}
\nomenclature[S]{\(\Delta f\)}{Resolução em frequência da DFT\nomunit{\si{\hertz}}}
\nomenclature[S]{\(T_s\)}{Duração de um símbolo LoRa\nomunit{\si{\second}}}

%% --- Modelo mecânico e térmico ---------------------------------------------
\nomenclature[M]{\(k\)}{Rigidez equivalente do caminho de transmissão\nomunit{\si{\newton\per\meter}}}
\nomenclature[M]{\(k_{\mathrm{base}}\)}{Rigidez da base, placa circular engastada\nomunit{\si{\newton\per\meter}}}
\nomenclature[M]{\(k_{\mathrm{res}}\)}{Rigidez do ressalto de centragem, compressão axial\nomunit{\si{\newton\per\meter}}}
\nomenclature[M]{\(k_{\mathrm{corpo}}\)}{Rigidez do corpo, viga tubular engastada\nomunit{\si{\newton\per\meter}}}
\nomenclature[M]{\(\mathcal{D}\)}{Rigidez flexional da placa\nomunit{\si{\newton\meter}}}
\nomenclature[M]{\(E\)}{Módulo de elasticidade\nomunit{\si{\pascal}}}
\nomenclature[M]{\(\nu\)}{Coeficiente de Poisson}
\nomenclature[M]{\(I\)}{Momento de inércia da seção tubular quadrada\nomunit{\si{\meter\tothe{4}}}}
\nomenclature[M]{\(A\)}{Área da seção resistente\nomunit{\si{\meter\squared}}}
\nomenclature[M]{\(L\)}{Comprimento do elo\nomunit{\si{\meter}}}
\nomenclature[M]{\(h\)}{Espessura da base\nomunit{\si{\meter}}}
\nomenclature[M]{\(a\)}{Raio da placa circular\nomunit{\si{\meter}}}
\nomenclature[M]{\(\ell\)}{Lado externo da seção tubular\nomunit{\si{\meter}}}
\nomenclature[M]{\(t\)}{Espessura da parede tubular\nomunit{\si{\meter}}}
\nomenclature[M]{\(f_n\)}{Frequência natural de montagem do conjunto\nomunit{\si{\hertz}}}
\nomenclature[M]{\(R_{\mathrm{th}}\)}{Resistência térmica do caminho condutivo\nomunit{\si{\kelvin\per\watt}}}

%% --- Descritores e avaliação ------------------------------------------------
\nomenclature[D]{\(x_{\mathrm{rms}}\)}{Valor eficaz do sinal na janela\nomunit{\si{\meter\per\second\squared}}}
\nomenclature[D]{\(\kappa\)}{Curtose do sinal na janela}
\nomenclature[D]{\(\sigma\)}{Desvio-padrão do sinal na janela\nomunit{\si{\meter\per\second\squared}}}
\nomenclature[D]{\(s(x,n)\)}{Escore de anomalia do Isolation Forest}
\nomenclature[D]{\(\mathrm{AUC}_\alpha\)}{Área parcial sob a curva ROC, acumulada até \(\alpha\)}
\nomenclature[D]{\(\alpha\)}{Limite superior da faixa de falso alarme na área parcial}
\nomenclature[D]{\(q\)}{Risco alvo na calibração de limiar por valores extremos}
\nomenclature[D]{\(N\)}{Número de amostras na janela}

%% --- Acrescentados pela fundamentação teórica ---------------------------------
\nomenclature[M]{\(r\)}{Razão entre frequência de excitação e frequência natural}
\nomenclature[M]{\(T(r)\)}{Transmissibilidade de base}
\nomenclature[M]{\(\zeta\)}{Razão de amortecimento}
\nomenclature[M]{\(R\)}{Resistência térmica\nomunit{\si{\kelvin\per\watt}}}
\nomenclature[D]{\(s(x,\psi)\)}{Escore de anomalia do Isolation Forest}
\nomenclature[D]{\(h(x)\)}{Comprimento do caminho de isolamento de \(x\)}
\nomenclature[D]{\(c(\psi)\)}{Comprimento médio de caminho para subamostra de \(\psi\) pontos}
\nomenclature[D]{\(\psi\)}{Tamanho da subamostra de cada árvore}
\nomenclature[D]{\(\beta\)}{Limite superior de falso alarme da área parcial}


\begin{document}

\selectlanguage{brazil}
\frenchspacing

%%%%%%%%%%%%%%%%%%%%%%%%%%%%%%%%%%%%%%%%%%%%%%%%%%%%%%%%%%%%%%%%%%%%%%%%%%%%%%%
%% ELEMENTOS PRÉ-TEXTUAIS
%%%%%%%%%%%%%%%%%%%%%%%%%%%%%%%%%%%%%%%%%%%%%%%%%%%%%%%%%%%%%%%%%%%%%%%%%%%%%%%

\imprimircapa
\imprimirfolhaderosto*

%% Ficha catalográfica: elaborada pela biblioteca da instituição; inserir o PDF
%% recebido aqui, se exigido.
%\begin{fichacatalografica}
%  \includepdf{fig_ficha_catalografica.pdf}
%\end{fichacatalografica}

\begin{folhadeaprovacao}
  \begin{center}
    {\ABNTEXchapterfont\large\imprimirautor}

    \vspace*{\fill}\vspace*{\fill}
    \begin{center}
      \ABNTEXchapterfont\bfseries\Large\imprimirtitulo
    \end{center}
    \vspace*{\fill}

    \hspace{.45\textwidth}
    \begin{minipage}{.5\textwidth}
        \imprimirpreambulo
    \end{minipage}%
    \vspace*{\fill}
   \end{center}

   Trabalho aprovado. \imprimirlocal, [dia] de [mês] de \imprimirdata:

   \assinatura{\textbf{\imprimirorientador} \\ Orientador}
   \assinatura{\textbf{[Nome do membro da banca]} \\ Convidado 1}
   \assinatura{\textbf{[Nome do membro da banca]} \\ Convidado 2}

   \begin{center}
    \vspace*{0.5cm}
    {\large\imprimirlocal}
    \par
    {\large\imprimirdata}
    \vspace*{1cm}
  \end{center}
\end{folhadeaprovacao}

%% Dedicatória, agradecimentos e epígrafe são opcionais (NBR 14724).
%\begin{agradecimentos}
%  ...
%\end{agradecimentos}

%%%%%%%%%%%%%%%%%%%%%%%%%%%%%%%%%%%%%%%%%%%%%%%%%%%%%%%%%%%%%%%%%%%%%%%%%%%%%%%
%% Resumo --- ABNT NBR 6028:2021
%%   parágrafo único; voz ativa; sem citações;
%%   palavras-chave separadas por ponto e vírgula, iniciais minúsculas,
%%   finalizadas por ponto.
%%
%% Enuncia a MESMA afirmação central do Capítulo 1 e do Capítulo 5, com a mesma
%% força. Ao alterar a tese, alterar os três.
%%%%%%%%%%%%%%%%%%%%%%%%%%%%%%%%%%%%%%%%%%%%%%%%%%%%%%%%%%%%%%%%%%%%%%%%%%%%%%%

\begin{resumo}
A manutenção preditiva de máquinas rotativas por análise de vibração é normalizada, e o
obstáculo à sua adoção é o custo de instrumentação, que restringe a cobertura aos ativos
críticos. Nós sensores de baixo custo respondem a esse obstáculo, mas a literatura que os
propõe raramente declara a banda e o piso de ruído que sustentam o desempenho reportado.
Este trabalho determina o que a cadeia de medição especificada para um dispositivo dessa
classe --- acelerômetro de consumo com taxa de amostragem de 1~kHz, microcontrolador de
baixo consumo, enlace de longo alcance e invólucro impresso em polímero --- permite medir,
e calibra o estágio de decisão de modo que o desempenho seja verificável no ponto de
operação pretendido. O método constrói condições de referência cujo resultado correto
decorre de solução analítica ou de definição normativa, submete-as à cadeia proposta e mede
o desvio. Os resultados têm três procedências, declaradas individualmente: sinal
sintetizado, dado real de bancada decimado em \emph{software} e modelagem sobre as
especificações de projeto. Em dado real, a 2\% de falso alarme, a detecção é de
75,0\% $\pm$ 3,5\% em desbalanceamento, entre 14,8\% e 20,7\% em defeito de rolamento puro
e entre 5,0\% e 7,2\% em desalinhamento. Em sinal sintetizado, três decisões de estimação
alteram os resultados por margens maiores que o efeito físico investigado: a integração no
domínio do tempo erra o indicador normativo em 73\% sem viés e em 716\% sob viés contido na
especificação do sensor, enquanto a integração no domínio da frequência o recupera
exatamente; a política de limiar padrão do detector marca 42,1\% da operação normal como
anomalia, contra 1,4\% sob limiar calibrado; e a partição aleatória de janelas infla a
acurácia em cerca de 36 pontos percentuais frente à partição por ensaio. Por modelagem, a
frequência de montagem do conjunto situa-se em 557~Hz, dentro da banda de medição, e o
limite térmico de operação é imposto pela bateria, e não pelo invólucro. Conclui-se que a
cadeia sustenta a detecção de desbalanceamento sob falso alarme fixado e não sustenta
diagnóstico de defeito de rolamento, e que a contribuição verificável de um dispositivo
dessa classe está no estágio de decisão, desde que o estágio de aquisição esteja declarado.
As conclusões valem para a cadeia especificada, e não para o desempenho do dispositivo
montado em operação industrial: a construção do protótipo e seu ensaio em ambiente real
constituem a etapa seguinte.

\vspace{\onelineskip}
\noindent
\textbf{Palavras-chave}: manutenção preditiva; análise de vibração; detecção de anomalia;
sistemas embarcados; instrumentação de baixo custo.
\end{resumo}

%%%%%%%%%%%%%%%%%%%%%%%%%%%%%%%%%%%%%%%%%%%%%%%%%%%%%%%%%%%%%%%%%%%%%%%%%%%%%%%
%% Abstract --- tradução do resumo, mesma afirmação e mesma força.
%%%%%%%%%%%%%%%%%%%%%%%%%%%%%%%%%%%%%%%%%%%%%%%%%%%%%%%%%%%%%%%%%%%%%%%%%%%%%%%

\begin{resumo}[Abstract]
\begin{otherlanguage*}{english}
Vibration-based predictive maintenance of rotating machinery is a standardised practice,
and the obstacle to its adoption is the cost of instrumentation, which restricts coverage
to critical assets. Low-cost sensor nodes address this obstacle, but the literature
proposing them rarely declares the measurement bandwidth and noise floor that underpin the
reported performance. This work determines what the measurement chain specified for a
device of this class --- a consumer-grade accelerometer sampling at 1~kHz, a low-power
microcontroller, a long-range radio link and a polymer printed enclosure --- is able to
measure, and calibrates the decision stage so that performance is verifiable at the
intended operating point. The method builds reference conditions whose correct result
follows from an analytical solution or from a normative definition, submits them to the
proposed chain and measures the deviation. Results have three provenances, declared
individually: synthesised signal, real test-rig data decimated in software, and modelling
over design specifications. On real data, at a 2\% false alarm rate, detection reaches
75.0\% $\pm$ 3.5\% for imbalance, between 14.8\% and 20.7\% for pure bearing faults and
between 5.0\% and 7.2\% for misalignment. On synthesised signals, three estimation
decisions alter the results by larger margins than the physical effect under
investigation: time-domain integration misses the normative indicator by 73\% without bias
and by 716\% under a bias within the sensor specification, whereas frequency-domain
integration recovers it exactly; the detector's default thresholding policy flags 42.1\% of
normal operation as anomalous, against 1.4\% under a calibrated threshold; and random
window partitioning inflates accuracy by roughly 36 percentage points relative to
partitioning by test run. From modelling, the assembly mounting frequency falls at 557~Hz,
inside the measurement band, and the thermal operating limit is imposed by the battery
rather than by the enclosure. We conclude that the chain supports imbalance detection under
a fixed false alarm rate and does not support bearing fault diagnosis, and that the
verifiable contribution of a device of this class lies in the decision stage, provided the
acquisition stage is declared. The conclusions hold for the specified chain, not for the
performance of the assembled device in industrial operation: building the prototype and
testing it in a real environment constitute the next stage.

\vspace{\onelineskip}
\noindent
\textbf{Keywords}: predictive maintenance; vibration analysis; anomaly detection;
embedded systems; low-cost instrumentation.
\end{otherlanguage*}
\end{resumo}


%% Listas
\pdfbookmark[0]{\listfigurename}{lof}
\listoffigures*
\cleardoublepage

\pdfbookmark[0]{\listtablename}{lot}
\listoftables*
\cleardoublepage

\printglossary[type=\acronymtype, title={Lista de abreviaturas e siglas}, toctitle={Lista de abreviaturas e siglas}]
\cleardoublepage

\printnomenclature
\cleardoublepage

\pdfbookmark[0]{\contentsname}{toc}
\tableofcontents*
\cleardoublepage

%%%%%%%%%%%%%%%%%%%%%%%%%%%%%%%%%%%%%%%%%%%%%%%%%%%%%%%%%%%%%%%%%%%%%%%%%%%%%%%
%% ELEMENTOS TEXTUAIS
%%%%%%%%%%%%%%%%%%%%%%%%%%%%%%%%%%%%%%%%%%%%%%%%%%%%%%%%%%%%%%%%%%%%%%%%%%%%%%%
\textual

%%%%%%%%%%%%%%%%%%%%%%%%%%%%%%%%%%%%%%%%%%%%%%%%%%%%%%%%%%%%%%%%%%%%%%%%%%%%%%%
%% Capítulo 1 --- Introdução
%%
%% Citações: abntex2cite (sistema autor-data). \cite{} → (AUTOR, ano);
%%           \citeonline{} → Autor (ano).
%%
%% MARCAÇÃO DE TRECHO EM ANDAMENTO --------------------------------------------
%% \emandamento{...} destaca em vermelho os trechos cuja validade muda quando o
%% protótipo for construído e ensaiado em ambiente real. É marcação de rascunho.
%%
%% Para REMOVER todas as marcas na versão final, troque a definição abaixo por:
%%     \newcommand{\emandamento}[1]{#1}
%% ou apague as chamadas \emandamento{...} do texto.
%%%%%%%%%%%%%%%%%%%%%%%%%%%%%%%%%%%%%%%%%%%%%%%%%%%%%%%%%%%%%%%%%%%%%%%%%%%%%%%

\makeatletter
\@ifundefined{emandamento}{%
  \@ifundefined{textcolor}%
    {\newcommand{\emandamento}[1]{#1}}%
    {\newcommand{\emandamento}[1]{\textcolor{red}{#1}}}%
}{}
\makeatother

\chapter{Introdução}
\label{cap:introducao}

A manutenção preditiva de máquinas rotativas por análise de vibração é uma prática
normalizada. A norma ISO 20816-3:2022 estabelece faixas de severidade a partir da
velocidade eficaz de vibração de banda larga e define os limites que separam operação
aceitável de operação inaceitável para máquinas industriais acima de
\SI{15}{\kilo\watt}, com rotação entre \num{120} e \num{30000}~r/min. Para essa medição,
a norma exige do instrumento resposta plana em faixa de \emph{ao menos}
\SIrange{10}{1000}{\hertz}~\cite{iso208163_2022}. O procedimento é
consolidado, e o obstáculo à sua adoção não é metodológico: é de custo de instrumentação.
Acelerômetros piezoelétricos com a banda exigida para diagnóstico, seus condicionadores e
os analisadores associados restringem a cobertura aos ativos considerados críticos, e
deixam sem instrumentação a maior parte do parque de motores de uma planta.

A resposta que a literatura recente oferece a esse obstáculo é o nó sensor de baixo custo:
uma unidade de medição inercial MEMS acoplada a um microcontrolador de baixo consumo, com
detecção de anomalia executada no próprio dispositivo e enlace de rádio de longo alcance
para telemetria. A viabilidade computacional dessa arquitetura está estabelecida. O
algoritmo Isolation Forest~\cite{liu2008} treina e executa dentro de um
ESP32~\cite{antonini2023}, existem variantes desenhadas para reduzir seu custo em
TinyML~\cite{barbariol2022}, e a combinação de microcontrolador, acelerômetro MEMS de
consumo e detecção não supervisionada já foi demonstrada em manutenção preditiva de baixo
custo~\cite{kolok2025}. O que a viabilidade computacional não resolve --- e é o ponto de
partida deste trabalho --- é a questão anterior: o que uma cadeia de medição dessa classe
efetivamente permite medir.

Essa questão tem uma ordem de precedência que organiza todo o restante. Primeiro vem a
aquisição: piso de ruído, banda passante e caminho mecânico entre a fonte de vibração e o
elemento sensor. Nenhuma decisão posterior recupera informação que o sensor não coletou.
Depois vem a estimação, que decide o que se faz com a informação disponível: a política de
limiar, a métrica de desempenho e o protocolo de validação. Por último vem a implantação
--- energia, comportamento térmico e invólucro ---, que determina por quanto tempo o
dispositivo continua operando. Os três elos são encadeados, e um erro no primeiro não é
compensável nos seguintes.

É contra essa cadeia que se identifica a lacuna. A literatura de vibração embarcada é
forte no primeiro elo e frágil no segundo. Sobre aquisição, há consenso quantitativo:
nenhum estudo de adequação de sensor revisado testou dispositivo com limite superior
abaixo de \SI{1500}{\hertz}~\cite{albarbar2008}, e os grupos que se propõem a diagnosticar
defeito de rolamento com MEMS escolhem sensores de dezenas de
quilohertz~\cite{landi2023,jakobsen2024}, porque a assinatura do defeito viaja modulada em
uma portadora de alta frequência~\cite{randall2011,antoni2007}. Sobre estimação, ao
contrário, o registro é lacunar. O escore do Isolation Forest tem \num{0,5} como
referência interpretativa, e não como corte calibrado~\cite{liu2008}; a implementação de
referência transporta essa interpretação para um deslocamento fixo, sem alterar o
ranqueamento~\cite{scikitlearn2026}; e os trabalhos que reportam bom desempenho com o
detector abandonam o padrão sem que a política adotada seja submetida a
varredura~\cite{antonini2023,brito2022}. Não se localizou, entre os trabalhos que executam
o detector em microcontrolador, nenhum que variasse o parâmetro de contaminação sobre a
mesma base e reportasse a fração de operação normal marcada como anomalia, nem taxa de
falso alarme publicada para detector treinado no próprio dispositivo --- que é justamente
o modo de operação exigido de uma máquina sem histórico rotulado. Soma-se a isso uma
fragilidade de protocolo com efeito documentado: janelas consecutivas de um mesmo ensaio
compartilham a assinatura da montagem, e sua divisão aleatória entre treino e teste
permite que o modelo identifique o ensaio em vez da falha. \citeonline{varejao2025}
denominam o fenômeno \emph{similarity bias} e relatam que, de \num{41} estudos revisados
sobre uma base de referência consolidada, \num{40} adotaram delineamento suscetível a ele.

Este trabalho ocupa essa lacuna com uma tese delimitada. Uma cadeia de medição de baixo
custo, com banda e piso de ruído declarados, sustenta a detecção de desbalanceamento sob taxa de falso
alarme fixada, mas não sustenta
diagnóstico de defeito de rolamento; e a contribuição verificável de um dispositivo dessa
classe está no segundo elo da cadeia --- limiar calibrado por risco alvo, métrica avaliada
no ponto de operação e protocolo de validação sem vazamento entre ensaios ---, contribuição
que só é válida se o primeiro elo estiver declarado em vez de disfarçado. O objeto de
estudo é o Kaelix, um dispositivo que combina acelerômetro MPU6050, microcontrolador
ESP32-S3, enlace LoRa e invólucro em ASA impresso por deposição, com detecção de anomalia
embarcada e rotulagem de severidade pela norma.

A validade dessa tese depende de que se declare a procedência de cada resultado, e ela não
é uniforme. Os resultados de detecção por modo de falha provêm de dado real: sinais de uma
bancada instrumentada, adquiridos a \SI{50}{\kilo\hertz} e submetidos a decimação
controlada em \emph{software} para emular a banda do dispositivo. O que é modelado, nesse
caso, não é o sinal, é a aquisição. Os resultados sobre banda, integração, limiar, métrica e
protocolo de validação provêm de sinal sintetizado, cuja resposta correta é conhecida, e
sustentam relações entre métodos, não desempenho em máquina. Os resultados de análise
mecânica e térmica provêm de modelagem computacional sobre as especificações de projeto,
tomadas como dadas. \emandamento{E a cadeia física do próprio
dispositivo não foi verificada: nenhum ensaio físico foi conduzido e nenhuma medição de
campo alimenta os modelos, uma vez que a construção do protótipo e seu ensaio em ambiente
real constituem a etapa seguinte deste trabalho.} Essa distinção restringe uma afirmação
específica e é declarada aqui por isso: \emandamento{os resultados sustentam conclusões
sobre o que a cadeia de aquisição especificada permite medir, e não sobre o desempenho do
dispositivo montado em operação industrial.} As duas coisas serão mantidas separadas ao
longo do texto, e cada resultado indicará de qual das três origens provém.

\section{Objetivo}
\label{sec:objetivo}

\subsection{Objetivo geral}
\label{ssec:objetivo-geral}

Determinar o que a cadeia de medição de vibração especificada para o Kaelix permite medir,
e calibrar o estágio de decisão do detector embarcado de modo que o desempenho reportado
seja verificável no ponto de operação pretendido.

\subsection{Objetivos específicos}
\label{ssec:objetivo-especifico}

Dentro desse objetivo geral, encontram-se como metas específicas:

\begin{itemize}
    \item Quantificar a perda de contraste do indicador de defeito de rolamento entre a
          cadeia de referência e a cadeia truncada do dispositivo, e verificar se o
          delineamento permite atribuir essa perda à banda ou ao método de demodulação;

    \item Determinar o método de integração de aceleração para velocidade compatível com a
          norma, e quantificar o erro dos métodos alternativos contra um sinal cuja
          integral é conhecida analiticamente~\cite{brandt2014,iso208163_2022};

    \item Medir a taxa de falso alarme produzida pela política de limiar padrão do
          Isolation Forest e o efeito do parâmetro de contaminação sobre ela;

    \item Comparar o desempenho do detector sob divisão aleatória por janela e sob divisão
          por ensaio, e quantificar a inflação de acurácia atribuível ao vazamento de
          assinatura de montagem;

    \item Avaliar, por modelagem, a resposta modal e a margem térmica do conjunto
          invólucro-sensor-fixação, e identificar qual componente impõe o limite de
          operação;

    \item Declarar, para cada questão levantada, se ela foi respondida por medição, por
          modelagem ou se permanece dependente de ensaio físico.
\end{itemize}

\section{Metodologia}
\label{sec:metodologia}

A estratégia adotada é comum às questões investigadas: construir uma condição de
referência cujo resultado correto decorre de solução analítica ou de definição normativa,
submetê-la à cadeia de medição proposta e medir o desvio introduzido pela configuração do
projeto. Quando a resposta certa é conhecida de antemão, um erro de implementação não
consegue se esconder atrás de um resultado plausível.

O trabalho se desenvolve em quatro etapas. A primeira é a revisão da literatura,
organizada por eixos correspondentes aos elos da cadeia descrita, com verificação de cada
referência em fonte primária e registro explícito das lacunas encontradas, incluindo
aquelas que a própria verificação derrubou. A segunda é a caracterização da aquisição: um sinal
sintetizado com resposta conhecida é submetido à cadeia de referência e à cadeia truncada, e
os descritores extraídos são comparados entre elas; em paralelo, sinais de bancada
adquiridos em alta taxa são decimados de forma controlada para emular a banda do
dispositivo e medir a detecção por modo de falha.
A terceira é a calibração do estágio de decisão, na qual a política de limiar é avaliada
por varredura, o desempenho é reportado no ponto de operação de baixa taxa de falso alarme
e o protocolo de validação é submetido à divisão por ensaio. A quarta é a verificação de
implantação, por modelagem mecânica e térmica do conjunto montado a partir das
especificações de projeto.

Duas decisões metodológicas atravessam as quatro etapas. A primeira é a rastreabilidade:
cada resultado é acompanhado da indicação de sua origem --- sinal sintetizado, dado real
decimado ou modelagem sobre especificação --- e os artefatos que o produzem, dados
intermediários e código de geração, são mantidos versionados e regeneráveis. A exceção são
os resultados de treino sobre dado real, que não têm artefato versionado próprio e exigem
nova execução do treino para serem reproduzidos (Seção~\ref{sec:dev-sinal}). A
segunda é o registro do que não foi verificado. As questões levantadas que dependem de
ensaio físico ou de consulta regulatória são enunciadas, mantidas abertas e associadas à
afirmação cujo alcance elas restringem, em vez de omitidas. Esse registro é parte do
método, e não uma ressalva sobre ele.

\section{Estrutura do Trabalho}
\label{sec:estrutura}

Este trabalho está organizado em seis capítulos. O Capítulo~\ref{cap:introducao}
apresentou o contexto, a lacuna identificada, a tese defendida e a procedência dos
resultados que a sustentam. O Capítulo~\ref{cap:fundamentacao} apresenta os conceitos
em que o trabalho se apoia: vibração de máquinas rotativas e suas assinaturas de falha,
aquisição e processamento de sinais, avaliação normativa de severidade, detecção de
anomalias e sua avaliação, e as restrições de um dispositivo embarcado. O
Capítulo~\ref{cap:revisao} revisa a literatura por eixos, seguindo a ordem
de precedência da cadeia de medição, e consolida o que a literatura sustenta, o que ela
contesta e quais lacunas permanecem abertas. O Capítulo~\ref{cap:desenvolvimento} descreve o desenvolvimento do
dispositivo: a cadeia de aquisição, o detector embarcado, o projeto de hardware e a
arquitetura de comunicação e recepção. O Capítulo~\ref{cap:resultados} apresenta os resultados e os confronta
com a literatura revisada, organizado pela ordem de precedência da cadeia: primeiro o que
a aquisição permite medir, depois o que a estimação decide com a informação disponível e,
por fim, o que a implantação impõe como limite de operação. O Capítulo~\ref{cap:conclusoes}
reúne as conclusões, delimita seu alcance e enuncia o trabalho futuro decorrente das
questões mantidas em aberto.

%%%%%%%%%%%%%%%%%%%%%%%%%%%%%%%%%%%%%%%%%%%%%%%%%%%%%%%%%%%%%%%%%%%%%%%%%%%%%%%
%% Capítulo 2 --- Fundamentação e revisão da literatura
%%
%% Reordenado pela ordem de precedência da cadeia, na MESMA sequência do
%% Capítulo 4: aquisição -> normalização -> estimação -> implantação.
%% Os relatórios de origem organizavam a revisão por "eixos" numerados 1, 3, 4 e
%% 5, abrindo pelo detector. A ordem aqui é a da cadeia, para que o capítulo de
%% revisão e o de resultados possam ser lidos em paralelo.
%%
%% A ausência de um eixo na numeração original era deliberada e está declarada
%% na Seção 2.1: as referências que o sustentariam não passaram na verificação.
%%%%%%%%%%%%%%%%%%%%%%%%%%%%%%%%%%%%%%%%%%%%%%%%%%%%%%%%%%%%%%%%%%%%%%%%%%%%%%%

\makeatletter
\@ifundefined{toprule}{%
  \newcommand{\toprule}{\hline}%
  \newcommand{\midrule}{\hline}%
  \newcommand{\bottomrule}{\hline}%
}{}
\@ifundefined{emandamento}{%
  \@ifundefined{textcolor}%
    {\newcommand{\emandamento}[1]{#1}}%
    {\newcommand{\emandamento}[1]{\textcolor{red}{#1}}}%
}{}
\makeatother

\chapter{Fundamentação e revisão da literatura}
\label{cap:revisao}

\section{Método da revisão}
\label{sec:rev-metodo}

A revisão foi conduzida por confronto, e não por levantamento: cada seção parte de uma
escolha de projeto ou de um resultado obtido e pergunta o que a literatura publicada diz
sobre ele --- se o confirma, se o contradiz ou se nada diz.

Todas as referências foram conferidas em fonte primária: depósito do editor, texto integral
do artigo, amostra oficial da norma ou folha de dados do fabricante. Onde o texto integral
não pôde ser acessado, isso consta na própria referência, e o achado correspondente é
tratado como de menor peso. Essa verificação não é formalidade: duas auditorias sucessivas
encontraram defeito em parcela substancial das referências inicialmente levantadas ---
primeiro autor trocado, iniciais incorretas colapsando autores distintos, norma retirada
citada como vigente e, em três casos, achado atribuído a uma fonte que não o continha.

Um eixo temático levantado na fase inicial não integra esta revisão. Nenhuma das
referências que o sustentariam sobreviveu à verificação em fonte primária, e a decisão foi
não citar o que não foi conferido. O registro fica aqui porque uma revisão que apresenta
apenas os eixos sobreviventes esconde o método que os selecionou.

A organização segue a ordem de precedência da cadeia de medição, a mesma do
Capítulo~\ref{cap:resultados}: primeiro a aquisição, que decide que informação existe;
depois a normalização, que impõe requisitos sobre o instrumento e sobre a estimativa; em
seguida a estimação, que decide o que se faz com a informação disponível; por fim a
implantação, que decide por quanto tempo o dispositivo continua operando.

\section{Aquisição: banda, ruído e detectabilidade de defeito de rolamento}
\label{sec:rev-aquisicao}

A tensão que organiza esta literatura é única: as frequências que se deseja medir em um
rolamento são baixas, mas o mecanismo físico que as torna mensuráveis é alto. O defeito
produz impactos periódicos na frequência característica da geometria, na ordem de dezenas
a centenas de hertz, e esses impactos excitam ressonâncias estruturais de alguns
quilohertz. A assinatura diagnóstica viaja modulada nessa portadora de alta frequência, e é
por demodulação de envelope sobre a banda de ressonância que ela é recuperada
\cite{randall2011,antoni2006,antoni2007}. O procedimento padrão localiza a banda ótima pela
díade que maximiza a impulsividade do sinal filtrado \cite{antoni2007}; um dispositivo cuja
frequência de Nyquist é de \SI{500}{\hertz} não tem díade a escolher.

Os estudos de adequação de sensor convergem para o limite superior de banda como a variável
determinante, e nenhum testa dispositivo próximo da faixa do Kaelix. A comparação de três
acelerômetros MEMS capacitivos contra um piezoelétrico de referência reprovou o de banda
mais estreita, de \SI{1500}{\hertz}, para monitoramento de condição \cite{albarbar2008}.
Os grupos que se propõem a diagnosticar rolamento com MEMS escolhem sensores de dezenas de
quilohertz \cite{landi2023,sharma2024}, e uma implementação de baixo custo publicada
resolve o problema por projeto, reservando um acelerômetro de banda estreita para
desbalanceamento e outro de \SI{40}{\kilo\hertz} para rolamento \cite{jakobsen2024}. Os
produtos comerciais medem, ao menos em um eixo, uma ordem de grandeza acima do teto de uma
IMU de consumo \cite{bentlynevada2025,erbessd2025,tractian2022}.

Há uma segunda objeção, anterior à banda e independente dela: o piso de ruído. A densidade
de ruído declarada em folha de dados para uma IMU de consumo, integrada sobre a banda
normativa, situa o sensor acima do ruído admissível para resolver a fronteira entre as
zonas normativas mais baixas. A literatura de MEMS levanta, de forma convergente, a
ausência de calibração rastreável desses dispositivos \cite{albarbar2008,rodrigues2021},
e a avaliação de sensibilidade diagnóstica de sensores digitais de catálogo chega à mesma
leitura \cite{fidali2024}. \emandamento{O orçamento de ruído usualmente atribuído à
literatura de fabricante não pôde ser confirmado em fonte primária \cite{looney2017}, e
fechar essa conta com densidade medida, e não declarada, permanece pendente.}

Um efeito adicional é pouco documentado e afeta diretamente dispositivos de consumo: a
ausência de filtro anti-\emph{aliasing} configurado abaixo da frequência de Nyquist. Sem
ele, conteúdo espectral acima da banda dobra para dentro da faixa de análise e torna-se
indistinguível de sinal legítimo. Não se localizou, na literatura revisada por pares de
diagnóstico de rolamento, descrição experimental do efeito desse artefato sobre descritores
estatísticos como curtose e fator de crista; ele aparece apenas em notas de aplicação de
fabricante, não conferidas nesta revisão.

\section{Normalização: o que a norma exige do instrumento e da estimativa}
\label{sec:rev-norma}

A avaliação de severidade em velocidade eficaz não é convenção arbitrária de unidades, mas
escolha justificada no próprio texto normativo \cite{iso208163_2022}. A norma vigente
estabelece as zonas de severidade para máquinas acima de \SI{15}{\kilo\watt} com rotação
entre \num{120} e \num{30000}~r/min, e exige do instrumento resposta plana em faixa de
\emph{ao menos} \SIrange{10}{1000}{\hertz}, com limite inferior de até \SI{2}{\hertz} em
máquinas próximas ou abaixo de \num{600}~r/min.

A distinção entre requisito de instrumento e janela de análise é decisiva e frequentemente
perdida. A exigência de resposta plana recai sobre o instrumento; é isso que torna o
truncamento de banda uma violação de requisito, e não uma escolha de processamento
reversível. A edição anterior da norma delega esse ``plana'' a uma norma de instrumentação
que o torna verificável, com tolerância de banda passante de aproximadamente
\SI{10}{\percent} entre \SI{12,6}{\hertz} e \SI{1000}{\hertz} e limitação de banda por par
de Butterworth de terceira ordem \cite{iso2954_2012}.

\emandamento{A edição de 2009 foi retirada e substituída pela de 2022
\cite{iso108163_2009,iso208163_2022}. Os valores de zona são idênticos nas duas, de modo
que a substituição afeta a rastreabilidade da citação e não os números.} A edição anterior
é mantida na bibliografia por ser a que a prática industrial e parte da literatura ainda
adotam, e por ser nela que está a delegação explícita à norma de instrumentação.

A conversão de aceleração para velocidade é obrigação normativa, e é também um problema
numérico com solução conhecida. A integração no domínio do tempo acumula qualquer
componente contínua presente no sinal, e o erro resultante cresce sem limite; a integração
no domínio da frequência, restrita à máscara normativa, é o método estabelecido, e a
comparação quantitativa entre os dois está publicada \cite{brandt2014}. O que não se
localizou publicado é a aplicação desse método ao indicador da norma em um dispositivo
truncado --- isto é, quanto um instrumento limitado a \SI{500}{\hertz} subestima o
indicador normativo frente à banda exigida.

\section{Estimação: detector, limiar, métrica e protocolo}
\label{sec:rev-estimacao}

Esta literatura desloca a pergunta de engenharia, mas o deslocamento é condicional, e a
condição precisa ficar à vista: \emph{dada banda de medição suficiente}, o que decide o
desempenho não é qual detector se escolhe, e sim como o limiar é calibrado, com que métrica
o resultado é lido e quanto do orçamento do microcontrolador o modelo consome. Sem banda,
nada disso importa --- nenhum corte de escore transforma contraste ausente em detecção.

\subsection{A escolha do detector não tem consenso}

Dois comparativos de grande porte discordam, e a discordância é informativa porque medem
coisas diferentes. O maior em cobertura de paradigmas avalia \num{30} algoritmos sobre
\num{57} conjuntos de dados e conclui que \emph{nenhum} método não supervisionado é
estatisticamente superior aos demais \cite{han2022}. A maior comparação restrita a métodos
não supervisionados, com \num{33} algoritmos sobre \num{52} conjuntos tabulares reais,
encontra vantagem para o Extended Isolation Forest, com a ressalva de que os dados se
separam em dois regimes \cite{bouman2024}. Não há contradição, e sim escopos distintos ---
mas há consequência: a evidência disponível recomenda uma variante que este trabalho não
adota.

A razão para permanecer no algoritmo canônico \cite{liu2008} é de custo de implementação
embarcada, e é verificável em estrutura: um nó da variante estendida armazena um vetor
normal com tantas componentes quantos forem os descritores, e a inferência passa de uma
comparação para um produto interno por nó. \emandamento{É argumento de projeto, e não
resultado de literatura; enquanto o orçamento embarcado não for medido no dispositivo, a
escolha permanece limitação de escopo declarada, não conclusão.}

O suporte específico para máquinas rotativas é, em contrapartida, o melhor que o algoritmo
canônico tem: entre \num{11} detectores não supervisionados, atinge F1 de \num{97,19} em
rolamento \emph{run-to-failure} e \num{97,20} em faltas mecânicas de bancada
\cite{brito2022}. Os dois primeiros conjuntos são adquiridos a \SI{20}{\kilo\hertz}, banda
que uma IMU de consumo não possui.

\subsection{O limiar padrão não é calibração}

No artigo fundador, o escore tem \num{0,5} como \emph{referência interpretativa}, e a
robustez demonstrada está nos parâmetros do \emph{ensemble}, não no corte \cite{liu2008}.
A implementação de referência explicita o mecanismo: o parâmetro de contaminação redefine
apenas um deslocamento fixo, nunca o escore subjacente \cite{scikitlearn2026}. Seguem daí
duas consequências mecânicas: o corte padrão é herdado de uma interpretação, e não estimado
nos dados; e toda métrica de ranqueamento, área sob a curva ROC inclusive, é invariante ao
parâmetro.

Os trabalhos que reportam bom desempenho abandonam o padrão --- fixando o escore
normalizado em \num{0,75} sob hipótese de treino limpo \cite{antonini2023}, ou derivando o
corte da razão de contaminação conhecida \cite{brito2022}. A alternativa principiada
dispensa hipótese sobre a distribuição do corpo dos dados e fixa o limiar a partir de um
risco alvo, que controla diretamente o número de falsos positivos \cite{siffer2017}.

O problema não é teórico. Três implementações independentes com o mesmo detector em nó
embarcado situam-se entre \num{22} e \SI{38}{\percent} de falso alarme sobre operação
normal \cite{chevtchenko2023,kolok2025}. Um dispositivo que marca um terço da operação
normal como anomalia não é utilizável em manutenção.

\subsection{A métrica e o protocolo de validação}

A área sob a curva ROC global não caracteriza o ponto de operação de um dispositivo de
campo, que precisa operar em falso alarme baixo. O benchmark de referência para o mesmo
regime de treino reporta área parcial restrita a essa faixa \cite{koizumi2020}, e a mesma
disciplina aparece em quem reporta PR-AUC além do F1 \cite{brito2022}.

Sobre o protocolo, a objeção é a mais séria e não foi identificada por este trabalho.
Janelas consecutivas de um mesmo ensaio compartilham a assinatura da montagem, e sua
divisão aleatória entre treino e teste permite que o modelo identifique o ensaio em vez da
falha. O fenômeno tem nome e quantificação: de \num{41} estudos revisados sobre uma base de
referência consolidada, \num{40} adotaram delineamento suscetível a ele \cite{varejao2025},
e a inflação medida por vazamento é da ordem de dezenas de pontos percentuais
\cite{vieira2026}. Trabalhos que revisam abordagens de classificação de falha de rolamento
chegam a leitura convergente \cite{abburi2023,smith2015}.

\section{Implantação: nós comparáveis, invólucro e energia}
\label{sec:rev-implantacao}

A literatura de nós sem fio para manutenção preditiva divide-se em duas populações que
raramente se citam. De um lado, produtos comerciais, com banda uma ordem de grandeza acima
da de uma IMU de consumo, caminho mecânico rígido até o elemento sensor e alimentação por
célula primária \cite{bentlynevada2025,erbessd2025,tractian2022}. De outro, nós acadêmicos
de baixo custo, que convergem para a mesma arquitetura estudada aqui: microcontrolador de
baixo consumo, IMU MEMS na ordem de \SI{1}{\kilo\hertz} e detecção não supervisionada
\cite{kolok2025,antonini2023,bisio2025}.

\subsection{O que a literatura de nós de baixo custo autoriza concluir}

O análogo mais próximo em hardware e em modo de falha reporta acurácia de
\SI{72,9}{\percent} $\pm$~\SI{2,4}{\percent}, com validação cruzada de cinco dobras sobre
cerca de \num{900} amostras e \emph{um único} modo de falha \cite{kolok2025}. O dado valida
a arquitetura como publicável. O que ele não faz é fixar teto de desempenho da classe: a
amostra é de um estudo único com um modo de falha, e o trabalho encerra em valor eficaz de
aceleração, sem integração para velocidade e sem critério normativo.

A comparação mais direta é embarcada e mede três detectores no mesmo nó e no mesmo
conjunto \cite{chevtchenko2023}. No mesmo hardware e no mesmo dado, um método baseado em
vizinhança domina o algoritmo adotado aqui tanto em especificidade quanto em latência,
sendo cerca de \num{26} vezes mais rápido. \emandamento{A resposta declarada é de custo de
memória --- o método de vizinhança exige manter o conjunto de referência, enquanto a
floresta armazena apenas limiares e índices --- mas o orçamento embarcado ainda não foi
medido, e a permanência no detector adotado é hipótese de projeto, não conclusão.}

O mesmo trabalho traz um achado que é restrição de implantação, e não detalhe: um motor com
cerca de um ano de uso tem assinatura normal distinta da de motores novos
\cite{chevtchenko2023}. Um modelo único com limiar fixo herda essa deriva do que conta como
normal entre máquinas e ao longo da vida --- cenário em que a calibração por risco alvo
deixa de ser refinamento e passa a ser requisito.

A viabilidade computacional, em contrapartida, está estabelecida: o detector treina e
executa dentro de um microcontrolador da classe \cite{antonini2023}, existe variante
desenhada para reduzir seu custo em TinyML \cite{barbariol2022}, e dispositivos da mesma
família comportam \emph{pipelines} consideravelmente mais caros \cite{uselis2025,kumar2026}.

\subsection{Invólucro e caminho mecânico}

É na física do dispositivo que a literatura contraria o projeto mais diretamente. A prática
consolidada em acelerometria exige que a frequência natural do conjunto montado seja
suficientemente superior à maior frequência de interesse, e os produtos comerciais obtêm
isso por base metálica aparafusada ou colada ao ativo. Nenhum dos nós de baixo custo
revisados publica a frequência natural do conjunto invólucro-sensor-fixação: um deles
reporta apenas autorressonâncias de placa e estrutura entre \SIrange{6}{9}{\kilo\hertz}
\cite{jakobsen2024}, outro adia explicitamente o invólucro \cite{antonini2023}, e para um
terceiro o dado não foi localizado na parte acessível do texto \cite{bisio2025}.

\subsection{Energia e margem térmica}

Na energia, a literatura contraria o projeto com clareza. O campo adota célula primária, e
portanto nunca carrega nada \cite{nabavi2025}; a revisão térmica disponível estabelece o
mecanismo de degradação por temperatura, mas não o limite de carga \cite{ma2018}. O limite
de carga de uma célula secundária vem da folha de dados e das normas de segurança
aplicáveis, não da literatura de manutenção preditiva. \emandamento{O orçamento térmico
completo --- carcaça quente, invólucro fechado e célula carregando --- não foi localizado
publicado, e é ele que decide se a escolha de bateria sobrevive.}

\section{Síntese}
\label{sec:rev-sintese}

Lidos em conjunto, os eixos não formam uma lista de pontos fortes e fracos: formam uma
cadeia com ordem de precedência. Primeiro a aquisição --- piso de ruído, banda e caminho
mecânico ---, e nenhuma decisão posterior recupera informação que o sensor não coletou.
Depois a estimação, que decide o que se faz com a informação disponível. Por último a
implantação, que decide por quanto tempo o dispositivo continua operando.

A conclusão desta revisão é que \textbf{a literatura é mais forte onde o dispositivo é mais
frágil, no primeiro elo, e mais frágil onde o dispositivo tem algo a dizer, no segundo.}

\subsection{O que a literatura sustenta}

O problema do limiar é real e não é preciosismo, e a alternativa principiada existe. A
disciplina de reportar a métrica no ponto de operação, e não apenas de forma agregada, tem
respaldo em benchmark de referência. A viabilidade computacional da arquitetura está
estabelecida, e a exigência de integrar aceleração para velocidade é normativa, com método
correto estabelecido. A forma do limite térmico --- margem imposta pela célula e não pelo
invólucro --- tem precedente comercial direto.

\subsection{O que a literatura contesta}

As objeções são de peso, e nenhuma se resolve com ajuste de \emph{software}. A banda:
nenhum estudo de adequação revisado testou sensor com limite superior próximo ao desta
classe, e o mecanismo físico da assinatura de rolamento não admite contorno. O piso de
ruído: a objeção é anterior à banda e não depende de nenhuma escolha de processamento. O
detector: a maior comparação restrita a métodos não supervisionados recomenda outra
variante, e no mesmo hardware embarcado um método de vizinhança mostrou-se mais específico
e mais rápido. \emandamento{Nenhuma dessas três objeções foi respondida por medição neste
trabalho; as respostas oferecidas são critérios de projeto declarados, e permanecem
hipóteses até que o orçamento embarcado seja medido.}

\section{Lacunas identificadas}
\label{sec:rev-lacunas}

Cada lacuna abaixo declara o que foi procurado, onde, e o que não foi encontrado. Duas
lacunas enunciadas na fase inicial caíram na verificação e estão registradas como caídas,
porque uma revisão que apresenta apenas as lacunas sobreviventes esconde o método.

\textbf{Caíram na verificação.} A afirmação de que não existiria falso alarme publicado
para o detector em nó embarcado é falsa: dois trabalhos o reportam ou permitem derivá-lo
\cite{chevtchenko2023,kolok2025}. A afirmação de que ninguém quantificaria a fração de
operação normal marcada pela política de limiar padrão também é falsa, pelo mesmo motivo.

\begin{enumerate}
  \item \textbf{Efeito isolado do parâmetro de contaminação.} Não se localizou trabalho que
        variasse o parâmetro sobre a \emph{mesma} base e o \emph{mesmo} modelo reportando a
        fração de operação normal marcada. O experimento que falta é barato e definitivo.
  \item \textbf{Falso alarme de detector treinado no próprio dispositivo.} O único trabalho
        revisado que treina \emph{on-device} é justamente o que declara a avaliação de
        acurácia fora de escopo \cite{antonini2023} --- e esse é o modo de operação exigido
        de uma máquina sem histórico rotulado.
  \item \textbf{Contraste diagnóstico como função contínua da banda.} A literatura discute
        o efeito da banda qualitativamente, e o único estudo revisado que faz varredura de
        taxa de amostragem o faz para acurácia de classificador supervisionado, não para
        contraste do indicador \cite{elbouharrouti2024}.
  \item \textbf{Artefato de \emph{aliasing} sobre descritores estatísticos.} Não localizado
        em literatura revisada por pares; aparece apenas em notas de aplicação de
        fabricante.
  \item \textbf{Custo do truncamento sobre o indicador normativo.} O erro dos métodos de
        integração está quantificado \cite{brandt2014}, mas não a aplicação ao indicador da
        norma em dispositivo truncado.
  \item \textbf{Detectabilidade limitada por piso de ruído.} Não se localizou análise de
        defeito mínimo detectável em função da densidade de ruído para IMUs de consumo.
  \item \textbf{Frequência natural do conjunto montado.} Nenhum nó de baixo custo revisado
        a publica.
  \item \textbf{Orçamento térmico de carga de célula secundária} em nó montado em carcaça
        de motor.
  \item \textbf{Separar efeito de banda de efeito de montagem.} Não se localizou benchmark
        aberto com sinais gravados \emph{simultaneamente} por MEMS de banda limitada e por
        piezoelétrico no mesmo ponto. Sem isso, os dois efeitos permanecem confundidos em
        qualquer medição de campo, inclusive nas deste projeto.
\end{enumerate}

As lacunas 1, 3 e 5 são as que o Capítulo~\ref{cap:resultados} aborda diretamente. As
demais permanecem abertas e reaparecem como trabalho futuro no
Capítulo~\ref{cap:conclusoes}.

%%%%%%%%%%%%%%%%%%%%%%%%%%%%%%%%%%%%%%%%%%%%%%%%%%%%%%%%%%%%%%%%%%%%%%%%%%%%%%%
%% Capítulo 3 --- Desenvolvimento e métodos
%%
%% Reúne o que os cinco relatórios de origem descreviam em aberturas de seção
%% separadas, cada um com escopo próprio. Os cinco "Escopo" colapsam aqui na
%% delimitação única da Seção 3.1.
%%
%% Ordem das subseções espelha a do Capítulo 4, para que método e resultado
%% correspondente possam ser localizados em paralelo.
%%%%%%%%%%%%%%%%%%%%%%%%%%%%%%%%%%%%%%%%%%%%%%%%%%%%%%%%%%%%%%%%%%%%%%%%%%%%%%%

\makeatletter
\@ifundefined{toprule}{%
  \newcommand{\toprule}{\hline}%
  \newcommand{\midrule}{\hline}%
  \newcommand{\bottomrule}{\hline}%
}{}
\@ifundefined{emandamento}{%
  \@ifundefined{textcolor}%
    {\newcommand{\emandamento}[1]{#1}}%
    {\newcommand{\emandamento}[1]{\textcolor{red}{#1}}}%
}{}
\makeatother

%% Ver nota sobre caminhos de figura no capitulo04.tex.
\providecommand{\kxfig}[1]{../../experiments/figures/output/#1}
\providecommand{\kxhw}[1]{../../#1}

\chapter{Desenvolvimento e métodos}
\label{cap:desenvolvimento}

\section{Delimitação}
\label{sec:dev-delimitacao}

O objeto de estudo é o Kaelix, um dispositivo de monitoramento de vibração de baixo custo
composto por acelerômetro MPU6050, microcontrolador ESP32-S3, enlace LoRa e invólucro em
ASA impresso por deposição, com detecção de anomalia embarcada e rotulagem de severidade
pela norma. As especificações de hardware, as cotas do desenho mecânico e a arquitetura de
processamento são tomadas como dadas; o que se determina é o que elas permitem medir.

A estratégia é comum a todas as questões investigadas: construir uma condição de referência
cujo resultado correto decorre de solução analítica ou de definição normativa, submetê-la à
cadeia proposta e medir o desvio introduzido pela configuração do projeto. Quando a
resposta certa é conhecida de antemão, um erro de implementação não consegue se esconder
atrás de um resultado plausível.

\emandamento{Nenhum ensaio físico foi conduzido e nenhuma medição de campo alimenta os
modelos. A construção do protótipo e seu ensaio em ambiente real são a etapa seguinte
declarada deste trabalho.} As três procedências de resultado que decorrem dessa delimitação
--- sinal sintetizado, dado real decimado e modelagem sobre especificação --- estão
definidas na Seção~\ref{sec:proc}, e cada valor do Capítulo~\ref{cap:resultados} é marcado
individualmente por elas.

\section{Modelo de sinal e cadeias de aquisição}
\label{sec:dev-sinal}

\subsection{Sinal sintetizado}

Os sinais foram sintetizados com dois mecanismos de escalas de frequência distintas.
Desbalanceamento e desalinhamento aparecem como componentes harmônicas da rotação, todas
abaixo de \SI{100}{\hertz}. Defeito de pista externa aparece como um trem de impulsos
amortecidos, repetido na taxa característica da geometria e modulando uma ressonância
estrutural. A aceleração simulada é

\begin{equation}
a(t) \;=\; \underbrace{\sum_{k=1}^{3} A_k \sin\!\left(2\pi k f_r t + \varphi_k\right)}_{\text{banda baixa}}
\;+\; \underbrace{\sum_{i} D\, e^{-\zeta_d (t-t_i)} \sin\!\left(2\pi f_{\mathrm{res}} (t-t_i)\right) u(t-t_i)}_{\text{impactos do defeito}}
\;+\; \varepsilon(t),
\label{eq:dev-sinal}
\end{equation}

\noindent
com $f_r$ a rotação, $A_k$ as amplitudes harmônicas, $D$ a amplitude do defeito, $\zeta_d$
a taxa de amortecimento do impacto, $f_{\mathrm{res}}$ a ressonância excitada,
$u(\cdot)$ o degrau unitário e $\varepsilon(t)$ ruído gaussiano branco. Os instantes de
impacto são $t_i = i/\mathrm{BPFO} + \delta_i$, em que $\delta_i$ modela o escorregamento
das esferas. As frequências características decorrem da geometria do rolamento pelas
expressões usuais de BPFO e BPFI.

\subsection{Três cadeias comparadas sobre o mesmo sinal}

A comparação isola o efeito da aquisição mantendo o sinal fixo:

\begin{enumerate}
  \item \textbf{Referência} --- amostragem a \SI{20}{\kilo\hertz}, representando um
        acelerômetro de banda larga.
  \item \textbf{Dispositivo corretamente configurado} --- decimação para
        \SI{1}{\kilo\hertz} com filtro anti-\emph{aliasing}.
  \item \textbf{Dispositivo sem filtro} --- subamostragem para \SI{1}{\kilo\hertz} sem
        filtro, representando a leitura direta do registrador quando o filtro passa-baixas
        digital interno não é configurado.
\end{enumerate}

No terceiro caso, uma componente em $f$ acima de Nyquist é observada na frequência aparente

\begin{equation}
f_{\mathrm{ap}} = \left| f - f_s \left\lfloor \frac{f}{f_s} + \frac{1}{2} \right\rceil \right|,
\label{eq:dev-alias}
\end{equation}

\noindent
com $f_s$ a taxa de amostragem e $\lfloor\cdot\rceil$ o arredondamento ao inteiro mais
próximo. É essa cadeia que corresponde ao estado atual do firmware
(Seção~\ref{sec:dev-firmware}).

\begin{figure}[!htb]
\centering
\includegraphics[width=0.88\textwidth]{\kxfig{fig8_circuito.pdf}}
\caption[Circuito de condicionamento e aquisição]{Circuito de condicionamento e aquisição. Os valores da figura são gerados
diretamente a partir das constantes do firmware, com o arquivo e a linha de origem
registrados no cabeçalho do gerador.}
\label{fig:circuito}
\end{figure}

Na comparação entre cadeias, os descritores extraídos por janela são o valor eficaz, a
curtose, o fator de crista, o pico e o pico a pico. O detector embarcado usa um conjunto
menor, idêntico no treino e no firmware: valor eficaz, curtose, fator de crista e frequência
dominante. A escolha tem consequência analisada no Capítulo~\ref{cap:resultados}: os
descritores descrevem energia, impulsividade e a linha espectral mais intensa, e nenhum
captura a relação entre harmônicas ou a fase.

\subsection{Dado real decimado}

Para as taxas de detecção por modo de falha, o sinal não é sintetizado. Utilizaram-se as
gravações do conjunto público MAFAULDA, obtidas em bancada instrumentada a
\SI{50}{\kilo\hertz} e decimadas em \emph{software} para \SI{1}{\kilo\hertz}, a banda do
dispositivo. %% TODO: citar o MAFAULDA; a chave ainda não existe em referencias.bib.
O conjunto soma \num{1951} arquivos em seis classes --- operação normal, desbalanceamento,
desalinhamento horizontal e vertical e defeito de rolamento nas posições \emph{overhang} e
\emph{underhang} ---, segmentados em \num{17559} janelas de \num{512} amostras. Apenas
\num{49} arquivos são de operação normal, e é deles que sai todo o treino, já que o detector
aprende somente a operação normal. A bancada está abaixo do escopo de potência e altura de
eixo da norma, o que a torna adequada como fonte de sinal para estudar a cadeia de medição e
inadequada como caso de aplicação normativa.

O limiar do detector é calibrado pelo quantil do escore sobre operação normal, com falso
alarme alvo de \SI{1}{\percent}. \emandamento{Os resultados desse treino não têm artefato
versionado próprio --- arquivo de resultados ou registro de execução ---: estão registrados
no histórico de alterações do repositório, e reproduzi-los exige executar o treino de novo.}

A validação emprega partição por ensaio, e não por janela, em cinco partições com os
arquivos de origem como grupos. A justificativa é o próprio
resultado de vazamento apresentado na Seção~\ref{sec:res-estimacao}: janelas do mesmo
ensaio compartilham a assinatura da montagem, e dividi-las aleatoriamente mede a
capacidade de identificar o ensaio, não a falha. Nos ensaios de rolamento, os resultados
são adicionalmente estratificados pelo nível de massa de desbalanceamento sobreposta,
porque o agregado confunde os dois efeitos.

\section{Conversão normativa e estágio de decisão}
\label{sec:dev-estimacao}

A norma define as zonas de severidade em velocidade eficaz de banda larga, e o acelerômetro
entrega aceleração \cite{iso208163_2022}. A conversão foi testada contra uma senoide cuja
integral é conhecida analiticamente, comparando cinco variantes: integração no tempo com e
sem filtro passa-alta, integração no domínio da frequência restrita à máscara normativa, e
esta última e a integração no tempo sem filtro repetidas sob viés contínuo contido na
especificação do sensor. O método correto
está estabelecido na literatura \cite{brandt2014}; o que se mede aqui é o erro das
alternativas sobre um valor de referência exato.

O estágio de decisão emprega detecção não supervisionada por Isolation Forest
\cite{liu2008}, treinada exclusivamente com dados de operação normal --- o cenário de uma
máquina sem histórico rotulado. Três aspectos foram submetidos a variação controlada: a
política de limiar, pela varredura do parâmetro de contaminação contra a fração de operação
normal marcada; a métrica, pela comparação entre área sob a curva ROC global e área parcial
restrita à região de falso alarme baixo \cite{koizumi2020}; e o protocolo de validação,
pela comparação entre partição aleatória por janela e partição por ensaio.

\section{Projeto mecânico e térmico}
\label{sec:dev-mecanica}

O invólucro foi analisado a partir das cotas do desenho técnico, por fórmulas fechadas de
resistência dos materiais. A rigidez do caminho de transmissão entre a superfície do motor
e o sensor foi modelada como associação em série de três elos --- a base como placa
circular engastada com carga central, o ressalto de centragem em compressão axial e o corpo
como viga tubular engastada:

\begin{equation}
k_{\mathrm{base}} = \frac{16\pi \mathcal{D}}{a^2},
\qquad
k_{\mathrm{res}} = \frac{E A}{L},
\qquad
k_{\mathrm{corpo}} = \frac{3 E I}{L^3},
\qquad
\frac{1}{k} = \sum_i \frac{1}{k_i},
\label{eq:dev-rigidez}
\end{equation}

\noindent
com $\mathcal{D} = E h^3 / [12(1-\nu^2)]$ a rigidez flexional da placa e
$I = [\ell^4 - (\ell-2t)^4]/12$ o momento de inércia da seção tubular quadrada. A
frequência de montagem e a transmissibilidade de base decorrem do modelo de um grau de
liberdade.

A análise térmica considera o invólucro fechado em regime permanente, com o calor entrando
pela carcaça do motor e o autoaquecimento do circuito como fonte interna. O caminho
condutivo é modelado como resistência térmica em série, e o limite de operação é
determinado pelo primeiro componente a atingir sua temperatura máxima admissível.

\emandamento{As condições de contorno são idealizadas: a rigidez real depende de pré-carga
de parafusos, planicidade das faces impressas e orientação das camadas, e a análise térmica
não modela resistência de contato, convecção forçada pelo ventilador do motor nem radiação
de superfícies próximas.}

\section{Geração e verificação do hardware}
\label{sec:dev-hardware}

Nenhum arquivo derivado é editado à mão. O esquemático, a placa e o sólido são gerados por
\emph{script} a partir de uma fonte única, e cada etapa tem verificação automática
associada, conforme a Tabela~\ref{tab:dev-cadeia}.

\begin{table}[!htb]
\centering
\caption{Cadeia de geração do hardware e verificação automática de cada etapa.}
\label{tab:dev-cadeia}
\small
\begin{tabular}{@{}lll@{}}
\toprule
Etapa & Saída & Verificação automática \\
\midrule
Esquemático & \texttt{.kicad\_sch}, netclasses & ERC com aviso tratado como erro \\
Placa       & \texttt{.kicad\_pcb}             & DRC com zonas preenchidas; colisão \\
Sólido      & \texttt{.step} da placa e do invólucro & modelo por \emph{footprint}; interferência \\
Modal       & frequências e formas             & convergência de malha em três refinos \\
Térmica     & campo e transiente               & balanço fechado contra modelo concentrado \\
\bottomrule
\end{tabular}
\end{table}

\begin{figure}[!htb]
\centering
\begin{minipage}{0.24\textwidth}\centering
\includegraphics[width=\linewidth]{\kxhw{hardware/pcb/v1/render_top.png}}\\[1pt]
{\footnotesize v1, superior}
\end{minipage}\hfill
\begin{minipage}{0.24\textwidth}\centering
\includegraphics[width=\linewidth]{\kxhw{hardware/pcb/v1/render_bottom.png}}\\[1pt]
{\footnotesize v1, inferior}
\end{minipage}\hfill
\begin{minipage}{0.24\textwidth}\centering
\includegraphics[width=\linewidth]{\kxhw{hardware/pcb/v2/render_top.png}}\\[1pt]
{\footnotesize v2, superior}
\end{minipage}\hfill
\begin{minipage}{0.24\textwidth}\centering
\includegraphics[width=\linewidth]{\kxhw{hardware/pcb/v2/render_bottom.png}}\\[1pt]
{\footnotesize v2, inferior}
\end{minipage}
\caption[Placas de circuito impresso geradas, versões v1 e v2]{Placas geradas nas duas versões. Contorno retangular com quatro recortes de canto
para os apoios do invólucro; duas camadas, com plano de terra em ambas.}
\label{fig:placas}
\end{figure}

A regeneração a partir de diretório limpo reproduz os arquivos de esquemático, placa e
símbolos byte a byte nas duas versões; o roteador é determinístico. Essa propriedade é o
que torna a verificação significativa: um resultado que não se reproduz não distingue
correção de coincidência.

As regras de fabricação adotadas são traço de \SI{0,20}{\milli\meter}, folga de
\SI{0,15}{\milli\meter} e via de \SI{0,60}{\milli\meter} com furo de
\SI{0,30}{\milli\meter}, contra mínimos declarados de \SI{0,15}{\milli\meter} para traço e
folga. \emandamento{Os defeitos de geração encontrados e corrigidos, bem como os que
permanecem abertos, constam do Apêndice~\ref{ape:hardware}.}

\begin{figure}[!htb]
\centering
\includegraphics[width=0.94\textwidth]{\kxfig{fig9_pdn.pdf}}
\caption[Rede de distribuição de alimentação]{Rede de distribuição de alimentação: impedância vista do pino de cada carga,
composição da indutância de laço e comprimento roteado contra distância em linha reta.}
\label{fig:pdn}
\end{figure}

\section{Firmware e estado de implementação}
\label{sec:dev-firmware}

O firmware é escrito em C++ para o ESP32-S3, com disciplina de sistemas críticos: caminhos
não implementados devolvem um estado de erro explícito e propagam-no, em vez de retornar
sucesso ou falha silenciosa. A justificativa dessa escolha e sua correspondência com a
prática de sistemas de alta criticidade constam do Apêndice~\ref{ape:praticas}.

\emandamento{Dois caminhos permanecem não implementados e são decisivos para a validade da
medição em campo: a inicialização do acelerômetro --- que inclui a conferência do registro
de identidade, o fundo de escala, o divisor de taxa de amostragem e, decisivamente, a
configuração do filtro passa-baixas interno abaixo de $f_s/2$ --- e o segundo caminho de
aquisição. Enquanto o filtro não for configurado, o dispositivo opera na terceira cadeia da
Seção~\ref{sec:dev-sinal}, aquela em que todo conteúdo acima de \SI{500}{\hertz} dobra para
dentro da banda de análise.}

\section{Comunicação e recepção}
\label{sec:dev-comunicacao}

A cadeia vai do motor ao manutentor em dez elos. Os parâmetros do enlace LoRa foram fixados
explicitamente --- fator de espalhamento, largura de banda e taxa de codificação ---, em vez
de herdados do padrão da biblioteca. A razão é de rastreabilidade: parâmetros herdados
significam que uma atualização de biblioteca altera o consumo do produto sem que a mudança
seja detectável.

O receptor assume três responsabilidades que o dispositivo não pode ter: carimbar o tempo,
já que o contador de partida fornece ordem e não instante; detectar lacuna na sequência,
porque perda silenciosa é indistinguível de máquina parada; e decidir quando abrir alerta.
A regra de disparo é implementada no receptor, e não no dispositivo, porque alterá-la no
dispositivo exigiria regravar cada aparelho em campo. A verificação de integridade é
simétrica: o CRC do receptor é comparado contra o do firmware compilado, em \num{204} casos
de teste.

\begin{figure}[!htb]
\centering
\includegraphics[width=\textwidth]{\kxfig{fig6_cadeia.pdf}}
\caption[Cadeia do sistema e estado de implementação de cada elo]{Cadeia do sistema e estado de implementação de cada elo. \emandamento{``Testado''
refere-se a verificação no \emph{host}: nenhum elo foi exercitado com rádio ou sensor
reais.}}
\label{fig:cadeia}
\end{figure}

\begin{figure}[!htb]
\centering
\includegraphics[width=0.92\textwidth]{\kxfig{fig5_topologia.pdf}}
\caption[Comparação entre arquiteturas topológicas]{Comparação entre arquiteturas topológicas. O enlace decide autonomia tanto quanto
decide alcance.}
\label{fig:topologia}
\end{figure}

\emandamento{A cadeia foi exercitada injetando os efeitos que o meio impõe --- colisão,
perda e corrupção de bit. O que a simulação não cobre é a propagação: alcance e margem de
enlace dependem de medição em campo, e nenhum elo foi exercitado com rádio ou sensor
reais.}

%%%%%%%%%%%%%%%%%%%%%%%%%%%%%%%%%%%%%%%%%%%%%%%%%%%%%%%%%%%%%%%%%%%%%%%%%%%%%%%
%% Capítulo 4 --- Resultados e discussão
%%
%% Organizado pela ordem de precedência da cadeia, como o Capítulo 1 promete:
%% aquisição -> estimação -> implantação, e então o confronto com a literatura.
%%
%% MARCAÇÃO DE PROCEDÊNCIA POR NÚMERO -----------------------------------------
%% Requisito estrutural deste capítulo, não recomendação de estilo.
%%
%% Os cinco relatórios de origem têm, cada um, procedência homogênea, declarada
%% uma vez no escopo --- exceto a revisão de literatura, que mistura as três. Os
%% três defeitos de procedência encontrados na auditoria estavam todos nela, e
%% em nenhum dos outros quatro. Este capítulo reproduz a condição de risco: ele
%% mistura as três procedências em texto corrido. Por isso a marcação é POR
%% NÚMERO (\pS, \pM, \pE) e não por seção.
%%
%% \pS  sinal sintetizado a partir do modelo da Seção 3.2
%% \pM  dado real de bancada instrumentada, decimado em software
%% \pE  modelagem ou cálculo sobre as especificações de projeto
%%
%%%%%%%%%%%%%%%%%%%%%%%%%%%%%%%%%%%%%%%%%%%%%%%%%%%%%%%%%%%%%%%%%%%%%%%%%%%%%%%

\makeatletter
%% booktabs não é carregado pelo template; se ausente, as réguas viram \hline
\@ifundefined{toprule}{%
  \newcommand{\toprule}{\hline}%
  \newcommand{\midrule}{\hline}%
  \newcommand{\bottomrule}{\hline}%
}{}
\@ifundefined{pS}{%
  \newcommand{\pS}{\textsuperscript{\textsf{S}}}%
  \newcommand{\pM}{\textsuperscript{\textsf{M}}}%
  \newcommand{\pE}{\textsuperscript{\textsf{E}}}%
}{}
\@ifundefined{emandamento}{%
  \@ifundefined{textcolor}%
    {\newcommand{\emandamento}[1]{#1}}%
    {\newcommand{\emandamento}[1]{\textcolor{red}{#1}}}%
}{}
\makeatother

%% Caminhos para as figuras geradas pelo pipeline do repositório.
%% NÃO copiar as figuras para dentro do TCC: elas são regeneráveis por
%% experiments/figures/ e uma segunda cópia divergiria em silêncio.
%% Ajustar o prefixo abaixo se a pasta de compilação mudar de lugar.
\providecommand{\kxfig}[1]{../../experiments/figures/output/#1}
\providecommand{\kxhw}[1]{../../#1}

\chapter{Resultados e discussão}
\label{cap:resultados}

\section{Procedência dos resultados}
\label{sec:proc}

Os resultados deste capítulo não têm origem uniforme, e a distinção decide o que cada um
sustenta. Três procedências se misturam no texto que segue, e cada valor numérico é
marcado individualmente:

\begin{description}
  \item[\pS\ Sinal sintetizado.] Gerado pelo modelo de dois mecanismos descrito na
        Seção~\ref{sec:dev-sinal}: harmônicas da rotação somadas a um trem de impulsos amortecidos que
        modula uma ressonância estrutural, mais ruído gaussiano. Sustenta afirmações sobre
        \emph{relações entre métodos} --- qual procedimento recupera o valor correto, qual
        decisão altera o resultado e em que direção. Não sustenta afirmação sobre
        desempenho em máquina.
  \item[\pM\ Dado real decimado.] Sinais de bancada instrumentada, adquiridos em alta taxa
        e decimados em \emph{software} para emular a banda do dispositivo. O que é
        modelado é a aquisição, não o sinal. Sustenta afirmações sobre \emph{taxas de
        detecção por modo de falha}.
  \item[\pE\ Modelagem sobre especificação.] Fórmulas fechadas e cálculo sobre as cotas do
        desenho e a configuração implementada, tomadas como dadas. Sustenta afirmações
        sobre \emph{limites impostos pelo projeto}, dentro das condições de contorno
        declaradas na Seção~\ref{sec:dev-mecanica}.
\end{description}

\emandamento{Nenhuma das três é medição no dispositivo montado. A marcação por número, e
não por seção, é requisito deste capítulo: valores das três procedências aparecem em
parágrafos vizinhos, e foi exatamente essa vizinhança que produziu, nos relatórios de
origem, três atribuições incorretas detectadas em auditoria.}

\section{O que a aquisição permite medir}
\label{sec:res-aquisicao}

\subsection{A banda do sensor descarta a evidência de defeito de rolamento}

O acelerômetro tem taxa máxima de amostragem de \SI{1}{\kilo\hertz}\pE, o que situa a
frequência de Nyquist em \SI{500}{\hertz}\pE. A frequência característica de defeito de
pista externa, \SI{104,6}{\hertz}\pE\ na geometria considerada, cai dentro dessa banda; a
ressonância estrutural de \SI{3,5}{\kilo\hertz}\pS\ que carrega a informação diagnóstica
não. O defeito repete dentro da banda, e a evidência do defeito não.

Sobre o sinal sintetizado, o contraste entre as condições com e sem defeito foi de
\num{413}\pS\ vezes por análise de envelope na banda de ressonância, amostrada a
\SI{20}{\kilo\hertz}, contra \num{20}\pS\ vezes por espectro direto a \SI{1}{\kilo\hertz}.
O par não isola o efeito da banda: entre as duas medidas mudam duas variáveis, a banda e o
método de demodulação, e dois pontos não constituem uma função. Ele sustenta que a cadeia
truncada perde contraste; não sustenta a atribuição quantitativa dessa perda à banda
isoladamente.

Um segundo efeito é de natureza distinta. Sem filtro passa-baixas configurado abaixo de
$f_s/2$, a ressonância de \SI{3,5}{\kilo\hertz}\pS\ não desaparece: reaparece rebatida em
\SI{500}{\hertz}\pS, frequência onde nada existe fisicamente. Descritores estatísticos
medidos sobre esse sinal podem separar as classes por um mecanismo alheio à física do
defeito --- uma medição ingênua pode funcionar pelos motivos errados e passar em validação.

\begin{figure}[!htb]
\centering
\includegraphics[width=\textwidth]{\kxfig{fig1_banda_sensor.pdf}}
\caption[O limite de banda descarta a evidência de defeito de rolamento]{\textbf{O limite de banda descarta a evidência de defeito de rolamento} (\pS).
\textbf{a}, Densidade espectral de potência para rolamento sadio e com defeito de pista
externa. A faixa sombreada marca a banda acessível ao dispositivo; os impactos excitam uma
ressonância inteiramente fora dela. \textbf{b}, Razão entre as amplitudes na frequência
característica nas duas condições. \textbf{c}, Razão entre as condições para quatro
descritores de impacto; nenhum sobrevive à limitação de banda com antialiasing correto.
\textbf{d}, Frequência aparente da ressonância após subamostragem sem filtro: a energia
reaparece onde nada existe fisicamente.}
\label{fig:banda}
\end{figure}

\subsection{O que o dispositivo detecta, por modo de falha}

Em dado real decimado, com o limiar calibrado para falso alarme alvo de \SI{1}{\percent},
o falso alarme medido foi de $\num{2,0} \pm \num{1,3}\,\si{\percent}$\pM. Nesse ponto de
operação, o quadro de detecção por classe é o da Tabela~\ref{tab:deteccao}.

\begin{table}[!htb]
\centering
\caption{Detecção por modo de falha a \SI{2}{\percent} de falso alarme medido, em dado
real decimado (\pM). Desvio entre partições onde registrado. As duas classes de rolamento
correspondem à condição sem massa de desbalanceamento sobreposta.}
\label{tab:deteccao}
\small
\begin{tabular}{lc}
\toprule
Modo de falha & Detecção \\
\midrule
Desbalanceamento                      & $\num{75,0} \pm \num{3,5}\,\si{\percent}$ \\
Rolamento \emph{overhang} (puro)      & \SI{20,7}{\percent} \\
Rolamento \emph{underhang} (puro)     & \SI{14,8}{\percent} \\
Desalinhamento vertical               & $\num{7,2} \pm \num{4,3}\,\si{\percent}$ \\
Desalinhamento horizontal             & $\num{5,0} \pm \num{3,1}\,\si{\percent}$ \\
\bottomrule
\end{tabular}
\end{table}

Apenas o desbalanceamento é detectado de forma útil. As demais classes ficam próximas do
piso, e duas observações sobre a tabela importam mais que os valores.

A primeira é sobre a dispersão. Nas duas classes de desalinhamento, o desvio entre
partições é da ordem do próprio valor: $\num{7,2} \pm \num{4,3}\,\si{\percent}$\pM\
admite intervalo que se aproxima de zero. Qualquer afirmação construída sobre esses dois
números é hipótese sustentada por evidência fraca, e não medição, e assim é tratada na
Seção~\ref{sec:res-discussao}.

A segunda é sobre estratificação. As classes de rolamento foram ensaiadas em quatro níveis
de massa de desbalanceamento sobreposta, e a detecção acompanha a massa, não o defeito:
sobe de \SI{14,8}{\percent}\pM\ sem massa para \SI{87,3}{\percent}\pM\ no maior nível, na
classe \emph{underhang}. Reportar o agregado das quatro condições --- \SI{42,6}{\percent}\pM\
--- como ``detecção de defeito de rolamento'' produziria exatamente o tipo de número
inflado que a Seção~\ref{sec:res-estimacao} documenta em outro contexto. O padrão indica
que, nos arquivos de rolamento, o modelo responde à massa adicionada, e não ao defeito.

\subsection{O invólucro filtra o sinal antes que ele alcance o sensor}

O invólucro integra a cadeia de medição, e não apenas a protege. A rigidez do caminho entre
a superfície do motor e o sensor resultou em \SI{1,45e6}{\newton\per\meter}\pE, o que para
uma massa de \SI{118}{\gram}\pE\ situa a frequência de montagem em \SI{557}{\hertz}\pE,
dentro da banda de medição. A transmissibilidade atinge aproximadamente cinco vezes em
\SI{500}{\hertz}\pE\ e o erro de amplitude ultrapassa \SI{10}{\percent} já em
\SI{168}{\hertz}\pE. A prática consolidada em acelerometria exigiria frequência de montagem
de ao menos \SI{3}{\kilo\hertz} para a banda normativa.

A decomposição da flexibilidade localiza o problema: a base responde por
\SI{51,1}{\percent}\pE\ e o corpo por \SI{48,5}{\percent}\pE, contra \SI{0,4}{\percent}\pE\
do ressalto. O resultado é pouco sensível à incerteza da espessura da base --- dobrando-a,
a frequência sobe apenas para \SI{750}{\hertz}\pE\ --- e a mesma geometria em alumínio
alcançaria \SI{3273}{\hertz}\pE. Duas hipóteses concorrentes foram testadas e descartadas:
a curvatura da carcaça, que reduz a retenção magnética entre \num{7} e \SI{30}{\percent}\pE\
preservando margem próxima de duas vezes a força inercial, e os modos próprios dos painéis,
situados acima de \SI{1300}{\hertz}\pE.

\begin{figure}[!htb]
\centering
\includegraphics[width=\textwidth]{\kxfig{fig3_involucro.pdf}}
\caption[O caminho elástico coloca a frequência de montagem dentro da banda]{\textbf{O caminho elástico coloca a frequência de montagem dentro da banda de
medição} (\pE). O caminho motor--base--ressalto--corpo--sensor resulta em frequência de
montagem de \SI{557}{\hertz}, com erro de amplitude acima de \SI{10}{\percent} a partir de
\SI{168}{\hertz}. A decomposição da flexibilidade mostra que base e corpo respondem por
praticamente toda ela. Os painéis correspondentes a fixação magnética e a modos de placa
são hipóteses testadas e descartadas.}
\label{fig:involucro}
\end{figure}

\begin{figure}[!htb]
\centering
\includegraphics[width=0.92\textwidth]{\kxfig{fig10_vibracao.pdf}}
\caption[Transmissibilidade do caminho de montagem e margem de fadiga]{Transmissibilidade do caminho de montagem, margem de fadiga de junta de solda por
componente e sensibilidade ao amortecimento e à severidade da excitação (\pE).}
\label{fig:vibracao}
\end{figure}

\begin{figure}[!htb]
\centering
\includegraphics[height=0.62\textheight]{\kxhw{hardware/fem/figuras/modos.png}}
\caption[Formas modais da placa e do caminho de montagem]{Formas modais da placa e do
caminho de montagem, obtidas por elementos finitos (\pE). A amplitude do autovetor é
arbitrária; a informação está na distribuição.}
\label{fig:modos}
\end{figure}

\section{O que a estimação decide}
\label{sec:res-estimacao}

Quatro decisões deste elo --- método de integração, política de limiar, métrica e
protocolo de validação --- alteraram os resultados por margens maiores que o efeito físico
sob investigação. Todos os valores desta seção são de sinal sintetizado, e o que eles sustentam
é a relação entre métodos, não o desempenho do dispositivo.

\subsection{A norma exige velocidade, e a integração precisa ser feita na frequência}

A norma define as zonas de severidade em velocidade eficaz, e o acelerômetro entrega
aceleração. Contra uma senoide de \SI{3,0}{\meter\per\second\squared} a \SI{50}{\hertz},
cuja integral é conhecida analiticamente e vale \SI{6,7524}{\milli\meter\per\second}, os
métodos comparam-se como na Tabela~\ref{tab:res-integracao}.

\begin{table}[!htb]
\centering
\caption{Velocidade eficaz por variante de integração, contra o valor analítico exato
(\pS).}
\label{tab:res-integracao}
\small
\begin{tabular}{lcc}
\toprule
Método & Resultado (\si{\milli\meter\per\second}) & Erro \\
\midrule
Integração no tempo, sem passa-alta          & \num{11,69}  & $+\SI{73}{\percent}$ \\
Integração no tempo, com passa-alta          & \num{17,31}  & $+\SI{156}{\percent}$ \\
\textbf{Integração na frequência, banda ISO} & \num{6,7525} & \textbf{\SI{0,0}{\percent}} \\
Na frequência, com viés de \SI{0,02}{\meter\per\second\squared} & \num{6,7525} & \SI{0,0}{\percent} \\
No tempo sem passa-alta, com o mesmo viés    & \num{55,08}  & $+\SI{716}{\percent}$ \\
\bottomrule
\end{tabular}
\end{table}

Dois resultados são contraintuitivos. O filtro passa-alta não corrige a integração no
tempo e chega a piorá-la, porque a soma cumulativa continua acumulando o transiente de
borda do filtro. E um viés de \SI{0,02}{\meter\per\second\squared}\pE, valor contido na
especificação do sensor, leva o resultado a \SI{55,08}{\milli\meter\per\second}\pS\ quando
o correto é \num{6,75}. O valor correto fica abaixo do limite entre as zonas C e D em três
das quatro combinações de grupo de máquina e tipo de fundação da norma; o valor obtido
excede esse limite em todas, e o sinal seria classificado na zona D qualquer que fosse a
máquina. A integração na frequência permanece exata sob o mesmo viés, porque o termo
contínuo cai fora da máscara imposta pela própria norma.

\subsection{O limiar padrão produz falso alarme incompatível com operação}

Treinado exclusivamente com dados de operação normal, o detector com o valor padrão do
parâmetro de contaminação classificou \SI{42,1}{\percent}\pS\ das janelas normais como
anomalia. Com o parâmetro fixado em \num{0,01}, o falso alarme caiu para
\SI{1,4}{\percent}\pS. A taxa de detecção permaneceu em \SI{100}{\percent}\pS\ e a área sob
a curva ROC em \num{1,000}\pS\ nos quatro valores testados.

Detecção unitária e área unitária são condições que nenhum conjunto real produz, e marcam
o alcance deste resultado: ele demonstra que o parâmetro desloca o corte sem alterar o
ranqueamento do detector, e é essa propriedade --- não o valor \SI{42}{\percent} --- que se
transporta para dado real.

\subsection{A métrica global não caracteriza o ponto de operação}

Sobre defeito incipiente, com descritores e modelo idênticos, a área sob a curva completa
caiu \SI{3,7}{\percent}\pS\ entre a cadeia de referência e a do dispositivo, enquanto a
área parcial restrita à região de falso alarme baixo caiu \SI{16,2}{\percent}\pS\ a
\SI{10}{\percent} de falso alarme e \SI{20,3}{\percent}\pS\ a \SI{5}{\percent}. A perda na
região em que um dispositivo de campo precisa operar é de \num{4,4} a \num{5,5} vezes a que
a área global indica.

\subsection{O protocolo de validação infla a acurácia}

Janelas consecutivas de um mesmo ensaio compartilham a assinatura da montagem --- nível de
ruído, ganho do sensor e condição de aperto. Divididas aleatoriamente entre treino e teste,
permitem que o modelo identifique o ensaio em vez da falha. O mesmo modelo, sobre os mesmos
dados e os mesmos descritores, atingiu \SI{99,0}{\percent}\pS\ de acurácia sob divisão
aleatória por janela e $\num{63,2} \pm \num{19,2}\,\si{\percent}$\pS\ sob divisão por
ensaio: inflação de aproximadamente \num{36} pontos percentuais.

A distribuição entre partições revela o mecanismo. Três das cinco ficaram próximas de
\num{0,50}\pS, o nível do acerto ao acaso, enquanto uma alcançou \num{1,00}\pS. O desvio
entre partições indica ausência de generalização --- informação que a média única do
protocolo incorreto suprime.

\begin{figure}[!htb]
\centering
\includegraphics[width=\textwidth]{\kxfig{fig2_metodologia.pdf}}
\caption[Quatro escolhas metodológicas alteram os resultados mais que o efeito medido]{\textbf{Quatro escolhas metodológicas alteram os resultados mais do que o efeito
medido} (\pS). \textbf{a}, Velocidade eficaz por cinco variantes de integração contra o
valor analítico exato; apenas a integração no domínio da frequência o recupera.
\textbf{b}, Deriva causada por viés contido na especificação do sensor. \textbf{c},
Acurácia do mesmo modelo sobre os mesmos dados, variando apenas a estratégia de partição;
pontos indicam as partições individuais. \textbf{d}, Fração da operação normal classificada
como anomalia por valor do parâmetro de contaminação. \textbf{e}, Curvas ROC por banda de
aquisição; a faixa sombreada marca a região de operação realista.}
\label{fig:metodologia}
\end{figure}

\section{O que a implantação impõe}
\label{sec:res-implantacao}

\subsection{A carga da bateria, e não o invólucro, define a margem térmica}

A hipótese inicial atribuía ao material do invólucro o papel de elo térmico mais frágil. A
análise inverteu essa leitura. A resistência térmica do caminho condutivo em ASA é de
\SI{88,8}{\kelvin\per\watt}\pE, contra \SI{0,074}{\kelvin\per\watt}\pE\ em alumínio, uma
razão de \num{1206}\pE\ vezes. Com base em alumínio, a temperatura interna alcançaria
\SI{79,0}{\celsius}\pE\ com a carcaça a \SI{80}{\celsius}. Com base em ASA, o balanço
térmico resolvido até a convergência dá \SI{33,0}{\celsius}\pE\ com a carcaça a
\SI{80}{\celsius}, e o modelo tridimensional do invólucro situa a placa entre
\SI{35,1}{\celsius}\pE\ e \SI{38,3}{\celsius}\pE\ com a carcaça a \SI{90}{\celsius}.

Em operação, nenhum limite é atingido até \SI{90}{\celsius} de carcaça: a placa permanece
abaixo tanto dos \SI{45}{\celsius} admitidos na carga da célula quanto dos
\SI{60}{\celsius} admitidos na descarga. O autoaquecimento é irrelevante em operação: a
potência média de \SI{2,0}{\milli\watt}\pE\ eleva a temperatura interna em
\SI{0,01}{\celsius}\pE, e todo o calor provém do exterior.

\emandamento{Durante a carga da bateria, não. O carregador linear da placa v3 dissipa cerca
de \SI{0,41}{\watt}\pE\ no pior ponto de uma carga a \SI{200}{\milli\ampere}, dentro do
invólucro fechado, e o modelo tridimensional eleva a placa em cerca de
\SI{15}{\kelvin}\pE\ por watt. O limite que restringe passa a ser o de carga da célula,
\SI{45}{\celsius}, e somente nessa fase: restam \SI{8,3}{\kelvin}\pE\ de folga com o
dispositivo fora do motor e de \SIrange{0,7}{3,9}{\kelvin}\pE\ com o motor a
\SI{90}{\celsius}. Esta última folga é otimista, porque o modelo aplica à troca entre a
placa e a cavidade o coeficiente de convecção da face externa, e a convecção dentro da
cavidade é mais fraca; a folga real pode ser nula. Na placa v3, o limite é imposto pelo
próprio carregador, que suspende a carga acima de \SI{45}{\celsius} e a retoma ao esfriar.
O procedimento recomendado é carregar o dispositivo fora do motor ou com o motor parado.}

Os demais limites ficam bem acima: circuitos integrados em \SI{85}{\celsius}, deformação
térmica do ASA em \SI{98}{\celsius} e ímãs em \SI{150}{\celsius}. A consequência de projeto
é um conflito entre elos: substituir o ASA por alumínio resolveria a limitação mecânica da
Seção~\ref{sec:res-aquisicao} e agravaria a térmica, levando a temperatura interna a
\SI{79,0}{\celsius}\pE\ sob a mesma carcaça, acima dos limites de carga e de descarga da
célula. Os dois requisitos recaem sobre a mesma peça e apontam em sentidos opostos.

\begin{figure}[!htb]
\centering
\includegraphics[width=\textwidth]{\kxfig{fig4_termica.pdf}}
\caption[O invólucro isola a eletrônica da carcaça]{\textbf{O invólucro isola a eletrônica
da carcaça} (\pE). Por ser mau condutor, o polímero do invólucro mantém a eletrônica abaixo
de todos os limites em operação. Os painéis \textbf{a} e \textbf{b} ainda apresentam o valor
do modelo concentrado não convergido, \SI{44,8}{\celsius} com a carcaça a
\SI{80}{\celsius}; o balanço convergido dá \SI{33,0}{\celsius}, e o modelo tridimensional
situa a placa entre \SI{35,1}{\celsius} e \SI{38,3}{\celsius} com a carcaça a
\SI{90}{\celsius}.}
\label{fig:termica}
\end{figure}

\begin{figure}[!htb]
\centering
\includegraphics[width=0.90\textwidth]{\kxfig{fig11_termica3d.pdf}}
\caption[Perfil de temperatura e transiente da parede da cavidade]{Perfil de temperatura ao longo da altura, transiente da parede da cavidade e
comparação com o modelo concentrado (\pE).}
\label{fig:termica3d}
\end{figure}

\begin{figure}[!htb]
\centering
\includegraphics[height=0.48\textheight]{\kxhw{hardware/fem/figuras/termica.png}}
\caption[Campo de temperatura em corte]{Campo de temperatura em corte (\pE). A placa não
integra o modelo de elementos finitos; sua temperatura provém de balanço radiativo
calculado separadamente.}
\label{fig:termicafem}
\end{figure}

\subsection{O enlace impõe um compromisso entre alcance e autonomia}

Os parâmetros do enlace foram fixados explicitamente, em vez de herdados do padrão da
biblioteca --- condição para que uma atualização não altere o consumo do produto sem que a
mudança seja detectável. O tempo no ar resultante é de \SI{185}{\milli\second}\pE, contra
os \SI{150}{\milli\second} que a estimativa inicial de consumo assumia sem derivação.

O custo em autonomia é direto: \SI{256}{\day}\pE\ no fator de espalhamento mais baixo
contra \SI{138}{\day}\pE\ no mais alto, com \SI{12,5}{\decibel} de ganho entre os extremos.
A comparação com Wi-Fi, modelada em \SI{5}{\second} a \SI{150}{\milli\ampere} por
transmissão, resulta em \SI{44}{\day}\pE\ contra \SI{236}{\day}\pE\ no enlace adotado.
Manter o rádio em \emph{standby} durante o repouso, em vez de colocá-lo em \emph{sleep},
resulta em \SI{45}{\day}\pE\ --- custo energético equivalente ao de adotar Wi-Fi.

\begin{figure}[!htb]
\centering
\includegraphics[width=0.92\textwidth]{\kxfig{fig7_gateway.pdf}}
\caption[Cadeia de comunicação exercitada por simulação de meio]{Cadeia de comunicação exercitada por simulação de meio (\pE), com injeção de
colisão, perda e corrupção de bit. \emandamento{A simulação não cobre propagação: alcance e
margem de enlace dependem de medição em campo.}}
\label{fig:gateway}
\end{figure}
\emandamento{A decisão definitiva depende da distância até o receptor, ainda não medida.}

\section{Confronto com a literatura}
\label{sec:res-discussao}

\subsection{Convergências}

O achado de inflação por vazamento de assinatura de montagem tem correlato direto: a
literatura mede inflação da ordem de dezenas de pontos percentuais em uma base de referência
consolidada, e relata que \num{40} de \num{41} estudos revisados sobre essa base adotaram
delineamento suscetível ao problema~\cite{varejao2025,vieira2026}. A convergência é de
direção e de ordem de grandeza: os \num{36} pontos\pS\ obtidos aqui vêm de sinal
sintetizado, e os valores publicados, de dado medido.

A fronteira de detecção encontrada reproduz, em dado real decimado, a divisão de trabalho que a
literatura implementou por projeto. Um nó publicado reserva um acelerômetro de banda
estreita para desbalanceamento e outro de dezenas de quilohertz para
rolamento~\cite{jakobsen2024}, e o trabalho mais próximo em arquitetura demonstra
desbalanceamento e adia rolamento~\cite{kolok2025}. Três caminhos independentes chegam à
mesma fronteira.

\subsection{Fragilidades desta análise}

O problema de calibração do limiar não é anomalia deste projeto: o escore do detector tem
\num{0,5} como referência interpretativa e não como corte calibrado~\cite{liu2008}, e a
implementação de referência transporta essa interpretação para um deslocamento
fixo~\cite{scikitlearn2026}. O resultado de \SI{42}{\percent}\pS\ obtido aqui vem de sinal
sintetizado, com detecção unitária, e por isso demonstra que a política padrão falha por
construção --- não que o dispositivo reproduza em máquina o falso alarme das implementações
publicadas, que se situam entre \num{22} e \SI{38}{\percent} sobre operação
normal~\cite{chevtchenko2023,kolok2025}. A comparação é de ordem de grandeza entre
conjuntos, máquinas e limiares diferentes.

A afirmação de que o desalinhamento falha por escolha de descritor, e não por banda, é a
observação mais original deste trabalho e a que se apoia na evidência mais frágil. O
mecanismo é claro: o desalinhamento se manifesta em harmônicas da rotação, muito abaixo de
\SI{500}{\hertz}, e portanto dentro da banda; a assinatura é a razão entre as linhas de
duas vezes e uma vez a rotação, que valor eficaz, curtose, fator de crista e frequência
dominante não capturam. Mas a detecção medida, $\num{7,2} \pm \num{4,3}\,\si{\percent}$\pM,
tem incerteza da ordem do próprio valor. A afirmação é mantida como hipótese acionável, com
o experimento que a testaria enunciado no Capítulo~\ref{cap:conclusoes}, e não como
resultado estabelecido.

\subsection{Escolhas de hardware contestadas}

A bateria recarregável contraria a prática estabelecida na literatura de nós de
monitoramento, que adota célula primária de longa duração~\cite{nabavi2025}. A escolha
introduz a restrição térmica de carga documentada na Seção~\ref{sec:res-implantacao}, que é
o limite de operação do dispositivo, e nenhum dos sistemas comparáveis a possui.

Sobre a banda, a objeção da literatura é direta e não admite contorno por
\emph{software}: nenhum estudo de adequação revisado testou sensor com limite superior
abaixo de \SI{1500}{\hertz}~\cite{albarbar2008}, e os grupos que diagnosticam rolamento com
MEMS escolhem sensores de dezenas de quilohertz~\cite{landi2023,jakobsen2024}, porque a
assinatura viaja modulada em portadora de alta frequência~\cite{randall2011,antoni2007}.
Os resultados desta seção são consistentes com essa objeção e a delimitam: a banda impede o
diagnóstico de rolamento e não impede a detecção de desbalanceamento.

%%%%%%%%%%%%%%%%%%%%%%%%%%%%%%%%%%%%%%%%%%%%%%%%%%%%%%%%%%%%%%%%%%%%%%%%%%%%%%%
%% Capítulo 5 --- Conclusões
%%
%% Par do capitulo01.tex: a Introdução estabelece a dívida, este capítulo a paga.
%% A afirmação central, a ordem de precedência da cadeia e a declaração de
%% procedência em três níveis são as mesmas do Capítulo 1, com a mesma força.
%%
%% Classe esperada: \documentclass[tcc,alf]{tcc-poli}
%% Citações: abntex2cite (sistema alfabético).
%% Pacotes usados: siunitx, hyperref (ambos já no principal.tex).
%%
%% MARCAÇÃO DE TRECHO EM ANDAMENTO --------------------------------------------
%% \emandamento{...} destaca em vermelho os trechos cuja validade muda quando o
%% protótipo for construído e ensaiado em ambiente real. Ver nota no capitulo01.
%%%%%%%%%%%%%%%%%%%%%%%%%%%%%%%%%%%%%%%%%%%%%%%%%%%%%%%%%%%%%%%%%%%%%%%%%%%%%%%

\makeatletter
\@ifundefined{emandamento}{%
  \@ifundefined{textcolor}%
    {\newcommand{\emandamento}[1]{#1}}%
    {\newcommand{\emandamento}[1]{\textcolor{red}{#1}}}%
}{}
\makeatother

\chapter{Conclusões}
\label{cap:conclusoes}

Este trabalho perguntou o que a cadeia de medição especificada para o Kaelix permite
medir, e não se um dispositivo de baixo custo pode ser construído. A resposta é que a
cadeia sustenta a detecção de desbalanceamento e de desvios de condição de operação sob
taxa de falso alarme fixada, e não sustenta diagnóstico de defeito de rolamento. O limite
não decorre de escolha de processamento e não é removível por ajuste de \emph{software}:
ele está na aquisição, e a ordem de precedência entre os elos da cadeia determina que
nenhuma decisão posterior recupere a informação que o sensor não coletou.

A contribuição verificável deste trabalho está no segundo elo. As decisões de limiar, de
métrica e de protocolo de validação alteraram os resultados por margens maiores que o
efeito físico sob investigação, e é sobre elas que os experimentos produziram números
novos. Essa contribuição depende de que o primeiro elo esteja declarado: um desempenho de
detecção reportado sem a banda e o piso de ruído que o produziram não é verificável, e a
prática de omiti-los é o que a revisão encontrou de mais frágil na literatura de vibração
embarcada.

\section{O que a aquisição permite medir}
\label{sec:conc-aquisicao}

A taxa máxima de amostragem de \SI{1}{\kilo\hertz} do acelerômetro situa a frequência de
Nyquist em \SI{500}{\hertz}. A frequência característica de defeito de pista externa, de
\SI{104,6}{\hertz} na geometria considerada, cai dentro dessa banda; a ressonância
estrutural de \SI{3,5}{\kilo\hertz} que carrega a informação diagnóstica não. Essa
separação é o resultado central do primeiro elo: a taxa de repetição do defeito estar
dentro da banda não significa que a evidência do defeito esteja.

Sobre o sinal simulado, o contraste entre as condições com e sem defeito foi de
\num{413} vezes por análise de envelope na banda de ressonância, amostrada a
\SI{20}{\kilo\hertz}, contra \num{20} vezes por espectro direto a \SI{1}{\kilo\hertz}.
Cabe a qualificação que a revisão já impõe a esse par: entre as duas medidas mudam duas
variáveis, a banda e o método de demodulação, e dois pontos não constituem uma função da
banda. O par sustenta a afirmação de que a cadeia truncada perde contraste; não sustenta a
atribuição quantitativa desse ganho à banda isoladamente. Separá-los exige varredura com
método fixo, e é trabalho futuro (Seção~\ref{sec:conc-futuro}).

Um segundo efeito é de natureza distinta e merece registro por ser, ele próprio, um
achado. Sem filtro passa-baixas configurado abaixo de $f_s/2$, a ressonância de
\SI{3,5}{\kilo\hertz} não desaparece: reaparece rebatida em \SI{500}{\hertz}, uma
frequência onde nada existe fisicamente. Descritores estatísticos medidos sobre esse sinal
podem separar as classes por um mecanismo que nada tem a ver com a física do defeito. Uma
medição ingênua pode, portanto, funcionar pelos motivos errados e passar em validação.

O invólucro integra a cadeia de medição, e não apenas a protege. A rigidez do caminho entre
a superfície do motor e o sensor situa a frequência de montagem do conjunto em
\SI{557}{\hertz}, dentro da banda de medição, com transmissibilidade de aproximadamente
cinco vezes em \SI{500}{\hertz} e erro de amplitude superior a \SI{10}{\percent} já em
\SI{168}{\hertz}. A prática consolidada em acelerometria exigiria frequência de montagem de
ao menos \SI{3}{\kilo\hertz} para a banda normativa. A decomposição da flexibilidade
localiza o problema na base e no corpo, responsáveis por \SI{51,1}{\percent} e
\SI{48,5}{\percent}, e não no ressalto, com \SI{0,4}{\percent}; a mesma geometria em
alumínio alcançaria \SI{3273}{\hertz}. Duas hipóteses concorrentes foram testadas e
descartadas: a curvatura da carcaça e os modos próprios dos painéis, estes situados acima
de \SI{1300}{\hertz}.

\section{O que a estimação decide}
\label{sec:conc-estimacao}

Três decisões do segundo elo alteraram os resultados por margens maiores que o efeito
medido.

A primeira é o método de integração. A norma define as zonas de severidade em velocidade
eficaz, e o acelerômetro entrega aceleração. Contra um sinal cuja integral é conhecida
analiticamente, a integração no domínio do tempo errou em $+\SI{73}{\percent}$ sem filtro
passa-alta e em $+\SI{156}{\percent}$ com ele, enquanto a integração no domínio da
frequência, restrita à máscara normativa, recuperou o valor exato. O filtro passa-alta não
corrige a integração no tempo e chega a piorá-la. Sob viés de
\SI{0,02}{\meter\per\second\squared}, valor contido na especificação do sensor, a
integração no tempo produziu \SI{55,08}{\milli\meter\per\second} onde o correto é
\num{6,75}: um sinal da zona A da norma seria classificado além da zona D. A integração na
frequência permanece exata sob o mesmo viés, porque o termo contínuo cai fora da máscara
imposta pela própria norma.

A segunda é a política de limiar. Treinado apenas com dados de operação normal, o detector
com o valor padrão do parâmetro de contaminação classificou \SI{42,1}{\percent} das
janelas normais como anomalia; com o parâmetro fixado em \num{0,01}, o falso alarme caiu
para \SI{1,4}{\percent}. A taxa de detecção permaneceu em \SI{100}{\percent} e a área sob
a curva ROC em \num{1,000} nos quatro valores testados. O parâmetro desloca o corte e não
altera a qualidade do detector, o que confirma que o desempenho reportado sem a política
de limiar declarada não é interpretável.

A terceira é a métrica. Sobre defeito incipiente, com descritores e modelo idênticos, a
área sob a curva completa caiu \SI{3,7}{\percent} entre a cadeia de referência e a do
dispositivo, enquanto a área parcial restrita à região de falso alarme baixo caiu
\SI{16,2}{\percent} a \SI{10}{\percent} de falso alarme e \SI{20,3}{\percent} a
\SI{5}{\percent}. A perda na região em que um dispositivo de campo precisa operar é quatro
a cinco vezes maior do que a área global indica.

A essas três soma-se o protocolo de validação. Janelas consecutivas de um mesmo ensaio
compartilham a assinatura da montagem, e sua divisão aleatória entre treino e teste permite
que o modelo identifique o ensaio em vez da falha. O mesmo modelo, sobre os mesmos dados e
os mesmos descritores, atingiu \SI{99,0}{\percent} de acurácia sob divisão aleatória por
janela e $\num{63,2} \pm \num{19,2}\,\si{\percent}$ sob divisão por ensaio, uma inflação de
aproximadamente \num{36} pontos percentuais. Três das cinco partições ficaram próximas de
\num{0,50}, o nível do acerto ao acaso, informação que a média única do protocolo
incorreto suprime.

\section{O que a implantação impõe}
\label{sec:conc-implantacao}

A hipótese inicial atribuía ao material do invólucro o papel de elo térmico mais frágil. A
análise inverteu essa leitura. A resistência térmica do caminho condutivo em ASA é de
\SI{88,8}{\kelvin\per\watt}, contra \SI{0,074}{\kelvin\per\watt} em alumínio, e com a
carcaça a \SI{80}{\celsius} a temperatura interna estabiliza em \SI{44,8}{\celsius}. O
autoaquecimento é irrelevante: a potência média de \SI{2,0}{\milli\watt} eleva a
temperatura interna em \SI{0,01}{\celsius}, de modo que todo o calor provém do exterior.

O componente que atinge primeiro seu limite é a bateria, pela restrição de carga em
\SI{45}{\celsius}, temperatura interna alcançada com carcaça em torno de
\SI{81}{\celsius}. O limite de operação do dispositivo é imposto pela célula, e não pelo
invólucro. A consequência de projeto é que substituir o material do invólucro por alumínio,
que resolveria a limitação mecânica da Seção~\ref{sec:conc-aquisicao}, agravaria a
limitação térmica, porque elevaria a temperatura interna a \SI{79,0}{\celsius} sob a mesma
carcaça. Os dois elos impõem requisitos opostos sobre a mesma peça.

\section{Alcance dos resultados}
\label{sec:conc-alcance}

As conclusões acima têm alcances distintos, porque a procedência dos resultados não é
uniforme, e enunciá-las sem essa distinção seria atribuir a todas a força da mais forte.

Os resultados de detecção em modos de falha provêm de dado real medido em bancada
instrumentada, submetido a decimação controlada em \emph{software} para emular a banda do
dispositivo: a \SI{2}{\percent} de falso alarme, \SI{75,0}{\percent} de detecção em
desbalanceamento, entre \num{15} e \SI{21}{\percent} em defeito de rolamento puro e entre
\num{5} e \SI{7}{\percent} em desalinhamento. O que é modelado nesses resultados é a
aquisição, não o sinal.

Os resultados de banda, de integração, de limiar e de protocolo provêm de sinais
sintetizados a partir de um modelo de dois mecanismos, e valem para a cadeia de
processamento, não para o desempenho em campo. Seus valores absolutos dependem dos
parâmetros de simulação --- amplitude do defeito, frequência de ressonância e nível de
ruído --- e devem ser lidos como ordens de grandeza. As relações entre métodos, que são o
objeto deste trabalho, não dependem desses parâmetros.

Os resultados mecânicos e térmicos provêm de modelagem analítica sobre as especificações
de projeto, tomadas como dadas. A análise mecânica emprega fórmulas fechadas com condições
de contorno idealizadas, e a rigidez real depende de pré-carga de parafusos, planicidade
das faces impressas e orientação das camadas. A análise térmica assume regime permanente e
não modela resistência de contato, convecção forçada pelo ventilador do motor nem radiação
de superfícies próximas.

\emandamento{A cadeia física do próprio dispositivo não foi verificada. Nenhum ensaio
físico foi conduzido e nenhuma medição de campo alimenta os modelos; a inicialização do
acelerômetro permanece não implementada no firmware e o filtro passa-baixas interno do
sensor é pendência declarada em código. A consequência é uma restrição específica, e não
uma ressalva genérica: os resultados sustentam conclusões sobre o que a cadeia de
aquisição especificada permite medir, e não sobre o desempenho do dispositivo montado em
operação industrial.} Em particular, a bancada sobre a qual os dados reais foram obtidos está abaixo
do escopo de potência e altura de eixo da norma, o que a torna adequada como fonte de
sinal para estudar a cadeia de medição e inadequada como caso de aplicação normativa.

\emandamento{Três dados externos permanecem pendentes e limitam o alcance das conclusões:
a temperatura real de carcaça dos equipamentos-alvo, a classe normativa dos motores a
monitorar e a validação dos carregadores dos conjuntos de dados públicos utilizados.}

\section{Trabalho futuro}
\label{sec:conc-futuro}

\emandamento{A construção do protótipo físico e seu ensaio em ambiente real constituem a
etapa seguinte declarada deste trabalho, e não uma possibilidade remota. Os resultados
apresentados aqui não substituem essa etapa: eles a preparam, ao estabelecer o que ela
precisa verificar e quais medidas tornam essa verificação interpretável. Um ensaio de
campo conduzido sem banda declarada, sem política de limiar explícita e sob divisão
aleatória de janelas reproduziria, em dado real, os mesmos artefatos que este trabalho
mediu em sinal controlado.}

Cada item abaixo é enunciado junto da afirmação que ele permitiria sustentar, e não como
lista de desejos.

\begin{itemize}
    \item \textbf{Varredura de taxa de amostragem com método fixo}, sobre a mesma gravação
          decimada em vários $f_s$, com decomposição entre efeito de banda e efeito de
          perda do envelope. É o que converte o par \num{413}$\times$/\num{20}$\times$ em
          função da banda, e o que falta para atribuir a perda de contraste à banda
          isoladamente.

    \item \textbf{Medição dos mesmos descritores com e sem filtro passa-baixas
          configurado}, no mesmo sinal. É o que estabelece experimentalmente o artefato de
          rebatimento sobre curtose e fator de crista, efeito para o qual não se localizou
          descrição na literatura revisada por pares.

    \item \textbf{Varredura do parâmetro de contaminação contra falso alarme medido}, em
          vibração industrial de banda limitada. É o que permite afirmar que a calibração
          por risco alvo supera a política padrão, em vez de apenas exibir que a política
          padrão falha.

    \item \textbf{Taxa de falso alarme de detector treinado no próprio microcontrolador.}
          É o modo de operação exigido de uma máquina sem histórico rotulado, e nenhum
          valor publicado foi localizado para ele.

    \item \textbf{Descritores de razão harmônica para desalinhamento.} O desempenho em
          desalinhamento, entre \num{5} e \SI{7}{\percent}, é o pior dos três modos
          medidos e não é explicável por banda: o modo se manifesta em harmônicas da
          rotação, muito abaixo de \SI{500}{\hertz}. A causa provável está nos descritores
          adotados, que capturam energia e impulsividade e não estrutura harmônica e de
          fase. É a pista mais acionável do conjunto, e não custa banda, sensor nem
          dinheiro.

    \item \textbf{Medição por excitação de impacto da frequência natural do conjunto
          montado}, incluindo a ressonância da fixação magnética. É o que substitui a
          estimativa analítica de \SI{557}{\hertz}, que depende de um amortecimento
          arbitrado, por valor medido.

    \item \emandamento{\textbf{Construção do protótipo e ensaio em ambiente real} ---
          etapa seguinte já prevista. É a única via que converte as conclusões sobre a
          cadeia especificada em conclusões sobre o dispositivo em operação. Três
          verificações a precedem no próprio dispositivo: implementar a inicialização do
          acelerômetro, configurar o filtro passa-baixas interno abaixo de $f_s/2$ e medir
          o piso de ruído efetivo, em vez de adotar o valor de folha de dados. Sem as duas
          primeiras, a medição de campo mede aliasing; sem a terceira, não há como separar
          defeito não detectado de defeito abaixo do piso.}
\end{itemize}

A ordem entre esses itens não é arbitrária. Os dois primeiros pertencem ao primeiro elo da
cadeia e condicionam a interpretação de todo o resto; os três seguintes pertencem ao
segundo elo, onde este trabalho localizou sua contribuição; o penúltimo pertence ao
terceiro. O último não é um item entre outros: é a continuação do trabalho, e o que este
documento oferece a ela é o conjunto de medidas que tornam seu resultado interpretável.

\emandamento{Nesse sentido, as conclusões aqui apresentadas são de uma etapa, e não de um
ciclo fechado. Elas delimitam o que a especificação permite esperar do dispositivo antes
que ele exista fisicamente, o que é precisamente o que permite projetar o ensaio de campo
de modo a testar afirmações, e não a confirmá-las.}


%%%%%%%%%%%%%%%%%%%%%%%%%%%%%%%%%%%%%%%%%%%%%%%%%%%%%%%%%%%%%%%%%%%%%%%%%%%%%%%
%% ELEMENTOS PÓS-TEXTUAIS
%%%%%%%%%%%%%%%%%%%%%%%%%%%%%%%%%%%%%%%%%%%%%%%%%%%%%%%%%%%%%%%%%%%%%%%%%%%%%%%
\postextual

\bibliography{../relatorio/referencias}

\begin{apendicesenv}
  \partapendices
  %%%%%%%%%%%%%%%%%%%%%%%%%%%%%%%%%%%%%%%%%%%%%%%%%%%%%%%%%%%%%%%%%%%%%%%%%%%%%%%
%% Apêndice A --- Registro de verificação do hardware
%%
%% Absorve o registro de engenharia do relatório de verificação de hardware:
%% correções no esquemático, defeitos de geração corrigidos e abertos, versões.
%% O material de PROJETO (cadeia de geração, placa, PDN) fica no Capítulo 3.
%%
%% PENDENTE DE TRANSCRIÇÃO: as tabelas completas de defeitos (itens 1--21) e de
%% comparação entre versões v1/v2 estão em docs/relatorio/verificacao-hardware.tex
%% e devem ser transportadas na íntegra. A estrutura abaixo é a do destino.
%%%%%%%%%%%%%%%%%%%%%%%%%%%%%%%%%%%%%%%%%%%%%%%%%%%%%%%%%%%%%%%%%%%%%%%%%%%%%%%

\chapter{Registro de verificação do hardware}
\label{ape:hardware}

Este apêndice reúne o registro de engenharia da cadeia de geração descrita na
Seção~\ref{sec:dev-hardware}. Ele é mantido no trabalho por uma razão de método: a
verificação só é informativa se os defeitos que ela encontrou forem publicados junto com os
que ela não encontrou.

Nenhum resultado aqui provém de bancada. Todos os valores são de geração, verificação
geométrica ou simulação, e o registro não reivindica validação experimental. Cada afirmação
tem associada a grandeza que a sustenta; itens sem medição associada aparecem como
\textsf{não verificado}.

\section{Correções no esquemático e defeitos corrigidos}
\label{sec:ape-corrigidos}

Os defeitos identificados e corrigidos durante a geração constam do relatório de
verificação com o valor antigo, o valor corrigido e a causa. Incluem, entre outros, erro de
referencial angular de \emph{pad} que produzia \num{42} curtos no DRC, e inversão de sinal
de rotação de \emph{footprint}.

\section{Defeitos abertos}
\label{sec:ape-abertos}

\emandamento{Os defeitos abaixo permanecem abertos e restringem afirmações específicas.}

\begin{itemize}
  \item \textbf{A placa não parafusa nos bosses do invólucro.} Na v1, o furo M3 exigiria
        material até \SI{28,2}{\milli\meter} na diagonal e o canto termina em
        \SI{27,5}{\milli\meter}; na v2, os recortes coincidem com os eixos dos bosses. A
        fixação passa a ser o assento pelo aro mais três furos M2 próprios.
  \item \textbf{Corredor do cabo.} Com \SI{6,0}{\milli\meter}, a v1 perde a saída de terra
        de dois componentes, com três itens soltos no DRC.
  \item \textbf{Altura do conjunto definida pelos conectores.} Decidir se o cabeçalho de
        gravação permanece populado no produto.
  \item \textbf{Divisor do NTC satura no frio}, atingindo \SI{3,22}{\volt} a
        \SI{-40}{\celsius}, acima da faixa útil do conversor. Abaixo de aproximadamente
        \SI{-35}{\celsius} os limiares de curto e de circuito aberto não disparam.
  \item \textbf{Numeração de pino do regulador varia por fabricante} ---
        \textsf{não verificado}; conferir contra a folha de dados antes de fabricar.
  \item \textbf{Furo de conector no invólucro sem conector correspondente} na netlist.
  \item \textbf{Rabicho de antena ausente da lista de materiais}; a ligação do conector
        do módulo de rádio ao painel não está especificada.
  \item \textbf{Nenhum código lê a tensão da bateria.} O caminho de hardware existe; os
        limiares do requisito de segurança correspondente não estão implementados.
\end{itemize}

\section{Limitações declaradas das análises}
\label{sec:ape-limitacoes}

A análise modal assume ASA isotrópico e maciço, hipótese otimista: a peça impressa com
preenchimento parcial é anisotrópica e menos rígida. O amortecimento do material não foi
medido, e sim herdado de análise anterior do projeto. As demais limitações constam da
Seção~\ref{sec:conc-alcance}.

  %%%%%%%%%%%%%%%%%%%%%%%%%%%%%%%%%%%%%%%%%%%%%%%%%%%%%%%%%%%%%%%%%%%%%%%%%%%%%%%
%% Apêndice B --- Práticas de sistemas críticos
%%
%% Absorve o relatório de práticas de sistemas críticos. Justifica escolhas de
%% engenharia de firmware; não toca a tese sobre a cadeia de medição.
%%
%% PENDENTE DE TRANSCRIÇÃO: a tabela de mecanismos por setor, com as colunas
%% Limite conhecido e Situação no Kaelix, está em
%% docs/relatorio/praticas-sistemas-criticos.tex e deve ser transportada.
%%%%%%%%%%%%%%%%%%%%%%%%%%%%%%%%%%%%%%%%%%%%%%%%%%%%%%%%%%%%%%%%%%%%%%%%%%%%%%%

\chapter{Práticas de sistemas críticos aplicadas ao firmware}
\label{ape:praticas}

Este apêndice registra a origem documentada das práticas de engenharia adotadas no
firmware. A unidade de análise é o \emph{mecanismo}, e não a norma: uma norma aparece como
origem de vários mecanismos.

Duas ressalvas delimitam o alcance. Primeira: aviônica civil e automotivo entram como
origem documentada de prática, não como alvo de certificação --- o dispositivo é protótipo
acadêmico e nenhum processo de certificação é reivindicado. Segunda: para cada mecanismo
registra-se também o que ele comprovadamente \emph{não} cobre; sem essa coluna, o
levantamento viraria lista de recomendações.

\section{Escalas de criticidade}
\label{sec:ape-escalas}

Não há equivalência normativa entre as escalas setoriais, e a correspondência citada na
literatura é informal. A norma de aviônica classifica software pela severidade da condição
de falha e não atribui taxa de falha ao software \cite{do178c}; a norma genérica é a única
das três com faixa numérica direta de medida de falha alvo \cite{iec61508}; a automotiva
compõe severidade, exposição e controlabilidade \cite{iso26262}.

\section{Mecanismos adotados}
\label{sec:ape-mecanismos}

Os mecanismos concentram-se em disciplina de codificação, detecção de erro e definição de
estado seguro --- categorias cujo custo é de projeto, e não de replicação de hardware, e
que portanto são acessíveis a um protótipo de baixo custo.

O mecanismo com consequência mais direta sobre este trabalho é a propagação explícita de
estado não implementado. Um caminho ainda não escrito devolve erro e propaga-o até o
receptor, em vez de retornar sucesso ou falhar silenciosamente. A justificativa é a mesma
que a literatura de sistemas críticos oferece para a redundância diversa: falha silenciosa
é indistinguível de operação correta, e é essa indistinguibilidade que remove a evidência
de que algo está errado \cite{knight1986,esa1996}.

\section{Limite conhecido das práticas adotadas}
\label{sec:ape-limites}

A disciplina de codificação e a propagação de erro não substituem redundância de hardware,
e a literatura documenta que a suposição de independência entre versões diversas de
software não se sustenta empiricamente \cite{knight1986}. \emandamento{Nenhuma das práticas
adotadas foi exercitada em hardware real; sua eficácia é verificada apenas em teste no
\emph{host}.}

\end{apendicesenv}

\end{document}
